% Step-by-step derivations of the mathematics in
%   G. Gryn'ova, M. L. Coote, C. Corminboeuf, "Theory and practice of uncommon molecular
%   electronic configurations", WIREs Comput. Mol. Sci. 2015, 5, 440-459. doi:10.1002/wcms.1233
% Every numbered result is tied to a verification script in ../verification/ (see the summary section).
% Build:  make derivations   (from the repository root)
\documentclass[11pt,a4paper]{article}
\usepackage[margin=2.3cm]{geometry}
\usepackage[T1]{fontenc}
\usepackage{lmodern,microtype}
\usepackage{amsmath,amssymb,amsthm,mathtools,bm}
\usepackage{booktabs,enumitem,xcolor}
\usepackage[most]{tcolorbox}
\usepackage[colorlinks=true,linkcolor=blue!50!black,citecolor=green!40!black,urlcolor=blue!60!black]{hyperref}
\usepackage[nameinlink]{cleveref}

\numberwithin{equation}{section}
\setlist{itemsep=2pt,topsep=3pt}

% ---- notation
\newcommand{\pa}{\varphi_a}
\newcommand{\pb}{\varphi_b}
\newcommand{\ket}[1]{\lvert #1\rangle}
\newcommand{\bra}[1]{\langle #1\rvert}
\newcommand{\braket}[2]{\langle #1\vert #2\rangle}
\newcommand{\mel}[3]{\langle #1\vert #2\vert #3\rangle}
\newcommand{\ev}[1]{\langle #1\rangle}
\newcommand{\eri}[4]{(#1#2\vert #3#4)}
\newcommand{\Hop}{\hat H}
\newcommand{\dd}{\,\mathrm d}
\newcommand{\abar}{\bar a}
\newcommand{\bbar}{\bar b}

% ---- boxes
\tcbset{colback=white,boxrule=0.5pt,arc=1pt,left=4pt,right=4pt,top=3pt,bottom=3pt,fonttitle=\bfseries\small}
\newtcolorbox{reviewclaim}{colframe=black!55,title=Statement in the review}
\newtcolorbox{resultbox}{colframe=green!45!black,title=Result}
\newtcolorbox{verifybox}{colframe=blue!55!black,title=Verification}
\newtcolorbox{sourcebox}{colframe=orange!70!black,title=Original sources (status)}
\newtcolorbox{caveatbox}{colframe=red!60!black,title=Caveat / correction}

\newtheorem{proposition}{Proposition}[section]

\title{Derivations for\\ \emph{``Theory and practice of uncommon molecular electronic configurations''}\\
\large (Gryn'ova, Coote \& Corminboeuf, \emph{WIREs Comput.\ Mol.\ Sci.} \textbf{2015}, 5, 440--459)}
\author{CH5103 term paper --- supporting derivations}
\date{\today}

\begin{document}
\maketitle

\begin{abstract}
The review contains a small number of explicit formulas and a large number of verbal claims whose
justification lives in the cited literature: Hund's rule and its competing explanations, static
versus dynamic correlation, disjoint non-bonding orbitals, spin polarisation, the Ovchinnikov rule,
the Janak/Koopmans view of singlet--triplet gaps, kinetic exchange, the arbitrariness of ROHF orbital
energies, the 3d/4s problem, and the critique of the aufbau principle. This document derives each
of them from first principles, one step at a time, states every assumption, and checks every
result symbolically (SymPy) or numerically (PySCF, exact diagonalisation). Where the original papers
could be read, their content is quoted; where they were paywalled, this is stated and only the
independently derived mathematics is given. Two errors in the review are identified: a sign error
in the two-determinant open-shell singlet (\cref{sec:twoe}), and the year of the Katriel--Pauncz
reference (1977, not 1975). One claim is shown to be system-dependent (same-spin vs.\ opposite-spin
correlation, \cref{sec:hund}).
\end{abstract}

\tableofcontents
\newpage

\setcounter{section}{-1}% sections are numbered 0-17, matching the file names and the docs
\section{Conventions}\label{sec:conv}

\begin{itemize}
  \item Atomic units throughout: $\hbar=m_e=e=4\pi\varepsilon_0=1$; energies in hartree ($E_h$).
  \item Non-relativistic, clamped-nuclei electronic Hamiltonian for $N$ electrons and nuclei $A$ with
        charges $Z_A$ at $\mathbf R_A$:
        \begin{equation}\label{eq:H}
          \Hop=\sum_{i=1}^{N}\hat h(i)+\sum_{i<j}\frac{1}{r_{ij}},\qquad
          \hat h(i)=-\tfrac12\nabla_i^2-\sum_A\frac{Z_A}{|\mathbf r_i-\mathbf R_A|}.
        \end{equation}
        $\Hop$ contains no spin operators.
  \item Spatial orbitals $\varphi_p(\mathbf r)$ are real and orthonormal, $\braket{\varphi_p}{\varphi_q}=\delta_{pq}$.
        Spin functions: $\alpha,\beta$, with $\braket{\alpha}{\alpha}=\braket{\beta}{\beta}=1$, $\braket{\alpha}{\beta}=0$.
        Spin orbitals: $\varphi_p\alpha\equiv p$, $\varphi_p\beta\equiv\bar p$.
        Shorthand: $\varphi_p(1)\equiv\varphi_p(\mathbf r_1)$, $\alpha_1\equiv\alpha(\sigma_1)$.
  \item One- and two-electron integrals (chemists' notation):
        \begin{align}
          h_{pq}&=\int\varphi_p(\mathbf r)\,\hat h\,\varphi_q(\mathbf r)\dd\mathbf r,\\
          \eri pqrs&=\iint\frac{\varphi_p(1)\varphi_q(1)\,\varphi_r(2)\varphi_s(2)}{r_{12}}\dd\mathbf r_1\dd\mathbf r_2,
        \end{align}
        with Coulomb and exchange integrals
        \begin{equation}\label{eq:JK}
          J_{pq}=\eri ppqq,\qquad K_{pq}=\eri pqpq .
        \end{equation}
        For real orbitals $\eri pqrs=\eri qprs=\eri pqsr=\eri rspq$ (8-fold symmetry).
  \item Unnormalised $2\times2$ Slater determinant of spin orbitals $\chi_1,\chi_2$:
        $D[\chi_1,\chi_2]=\chi_1(1)\chi_2(2)-\chi_2(1)\chi_1(2)$; the normalised determinant is
        $\tfrac1{\sqrt2}D$, written $\ket{\chi_1\chi_2}$.
  \item Tags used: \textbf{[paper]} = stated in the review; \textbf{[derived]} = derived here;
        \textbf{[verified: \dots]} = checked in this work against a source or a calculation.
\end{itemize}

\section{Antisymmetry and two-electron spin functions}\label{sec:anti}

\begin{reviewclaim}
``According to the Pauli exclusion principle, the wavefunction for a system of fermions must be
antisymmetric with respect to the interchange of any two particles \dots\ only the system (atom,
molecule) as a whole possesses well-defined (sharp) stationary states.'' (p.~441)
\end{reviewclaim}

\subsection{Why the exchange eigenvalue is $\pm1$}
Let $\hat P_{12}$ swap all coordinates (space and spin) of electrons 1 and 2:
$\hat P_{12}\Psi(x_1,x_2,\dots)=\Psi(x_2,x_1,\dots)$.
\begin{enumerate}
  \item Swapping twice restores the original ordering, so $\hat P_{12}^2=\hat 1$.
  \item If $\hat P_{12}\Psi=\lambda\Psi$, then $\Psi=\hat P_{12}^2\Psi=\lambda^2\Psi$, hence $\lambda=\pm1$.
  \item $\Hop$ of \cref{eq:H} is symmetric in the electron labels, so $[\Hop,\hat P_{12}]=0$, and the
        eigenstates of $\Hop$ can be chosen as simultaneous eigenstates of $\hat P_{12}$.
  \item The postulate (spin--statistics) for electrons selects $\lambda=-1$:
        \begin{equation}\label{eq:antisym}
          \Psi(x_2,x_1,\dots)=-\Psi(x_1,x_2,\dots).
        \end{equation}
  \item A Slater determinant obeys \cref{eq:antisym} automatically: exchanging electrons 1 and 2
        exchanges two rows of the determinant, which multiplies it by $-1$.
  \item If two spin orbitals coincide, two columns are equal and the determinant vanishes. This is
        Pauli's original, one-electron version (``no two electrons with the same four quantum
        numbers'') recovered as a \emph{consequence} of \cref{eq:antisym} in the independent-particle
        approximation. Section~\ref{sec:qn} shows why, beyond that approximation, one-electron quantum
        numbers do not exist.
\end{enumerate}

\subsection{Spin operators on two electrons}
Single-electron spin-$\tfrac12$ operators act as
$\hat s_z\alpha=\tfrac12\alpha$, $\hat s_z\beta=-\tfrac12\beta$, $\hat s_+\beta=\alpha$, $\hat s_+\alpha=0$,
$\hat s_-\alpha=\beta$, $\hat s_-\beta=0$.
Total spin: $\hat S_z=\hat s_z(1)+\hat s_z(2)$, $\hat S_\pm=\hat s_\pm(1)+\hat s_\pm(2)$ and
\begin{equation}\label{eq:S2}
  \hat S^2=\hat S_-\hat S_++\hat S_z^2+\hat S_z .
\end{equation}
(\emph{Proof of \cref{eq:S2}:} $\hat S_-\hat S_+=(\hat S_x-i\hat S_y)(\hat S_x+i\hat S_y)
=\hat S_x^2+\hat S_y^2+i[\hat S_x,\hat S_y]=\hat S^2-\hat S_z^2-\hat S_z$, using $[\hat S_x,\hat S_y]=i\hat S_z$.)

Apply \cref{eq:S2} to the four products $\alpha_1\alpha_2,\alpha_1\beta_2,\beta_1\alpha_2,\beta_1\beta_2$:
\begin{enumerate}
  \item $\alpha_1\alpha_2$: $\hat S_+\alpha_1\alpha_2=0$; $\hat S_z\alpha_1\alpha_2=\alpha_1\alpha_2$.
        Hence $\hat S^2\alpha_1\alpha_2=(0+1+1)\alpha_1\alpha_2=2\,\alpha_1\alpha_2$, i.e.\ $S(S+1)=2$, $S=1$, $M_S=1$.
  \item $\beta_1\beta_2$: by the same steps with $\hat S_+\beta_1\beta_2=\alpha_1\beta_2+\beta_1\alpha_2$,
        $\hat S_-(\alpha_1\beta_2+\beta_1\alpha_2)=2\beta_1\beta_2$, $\hat S_z^2+\hat S_z\to 1-1=0$:
        $\hat S^2\beta_1\beta_2=2\beta_1\beta_2$ ($S=1$, $M_S=-1$).
  \item $\alpha_1\beta_2$: $\hat S_+\alpha_1\beta_2=\alpha_1\alpha_2$, then $\hat S_-\alpha_1\alpha_2=\beta_1\alpha_2+\alpha_1\beta_2$;
        $\hat S_z\alpha_1\beta_2=0$. Hence $\hat S^2\alpha_1\beta_2=\alpha_1\beta_2+\beta_1\alpha_2$.
        Likewise $\hat S^2\beta_1\alpha_2=\alpha_1\beta_2+\beta_1\alpha_2$.
  \item Therefore $\hat S^2(\alpha_1\beta_2\pm\beta_1\alpha_2)=(1\pm1)(\alpha_1\beta_2\pm\beta_1\alpha_2)$:
\end{enumerate}
\begin{resultbox}
\vspace{-1.2em}
\begin{align}
  \Theta_{1,+1}&=\alpha_1\alpha_2,\quad
  \Theta_{1,0}=\tfrac1{\sqrt2}(\alpha_1\beta_2+\beta_1\alpha_2),\quad
  \Theta_{1,-1}=\beta_1\beta_2 &&(S=1,\ \text{symmetric under }1\leftrightarrow2),\label{eq:tripspin}\\
  \Theta_{0,0}&=\tfrac1{\sqrt2}(\alpha_1\beta_2-\beta_1\alpha_2) &&(S=0,\ \text{antisymmetric}).\label{eq:singspin}
\end{align}
Combined with \cref{eq:antisym}: a triplet must have an \emph{antisymmetric} spatial part, a singlet a
\emph{symmetric} spatial part.
\end{resultbox}
\begin{verifybox}
\texttt{t01\_spin\_and\_determinants.py}: $\hat S^2$, $\hat S_z$ built as $4\times4$ matrices; all four
eigen-relations confirmed exactly (SymPy). \textbf{PASS}.
\end{verifybox}

\section{Two electrons in two orbitals: the wavefunctions of Figure 1}\label{sec:twoe}

\begin{reviewclaim}
Triplet components (p.~441):
$^3\Psi(M_S=+1)=\tfrac1{\sqrt2}[\pa(1)\pb(2)-\pb(1)\pa(2)]\alpha_1\alpha_2$,
$^3\Psi(M_S=-1)=\tfrac1{\sqrt2}[\pa(1)\pb(2)-\pb(1)\pa(2)]\beta_1\beta_2$, ``i.e.\ they can be
represented as single determinants v and vi''. Open-shell singlet:
$^1\Psi(M_S=0)=\tfrac1{\sqrt2}[\pa(1)\pb(2)+\pb(1)\pa(2)]\cdot\tfrac1{\sqrt2}[\beta_1\alpha_2-\alpha_1\beta_2]$,
``a combination of two determinants, iii and iv'', printed as
$\tfrac1{\sqrt2}\tfrac1{\sqrt2}\bigl(\,|\dots|_{\text{iii}}+|\dots|_{\text{iv}}\bigr)$.
\end{reviewclaim}

\subsection{The six determinants}
With two spatial orbitals and two spin functions there are four spin orbitals $a,\abar,b,\bbar$ and
$\binom42=6$ two-electron determinants (Figure~1 of the review):
\begin{center}
\begin{tabular}{llll}
\toprule
label & occupation & $D$ (unnormalised) & $M_S$\\
\midrule
i   & $\pa^2$ & $D[a,\abar]=\pa(1)\pa(2)\,(\alpha_1\beta_2-\beta_1\alpha_2)$ & 0\\
ii  & $\pb^2$ & $D[b,\bbar]=\pb(1)\pb(2)\,(\alpha_1\beta_2-\beta_1\alpha_2)$ & 0\\
iii & $\pa^\alpha\pb^\beta$ & $D[a,\bbar]=\pa(1)\pb(2)\alpha_1\beta_2-\pb(1)\pa(2)\beta_1\alpha_2$ & 0\\
iv  & $\pa^\beta\pb^\alpha$ & $D[\abar,b]=\pa(1)\pb(2)\beta_1\alpha_2-\pb(1)\pa(2)\alpha_1\beta_2$ & 0\\
v   & $\pa^\alpha\pb^\alpha$ & $D[a,b]=[\pa(1)\pb(2)-\pb(1)\pa(2)]\alpha_1\alpha_2$ & $+1$\\
vi  & $\pa^\beta\pb^\beta$ & $D[\abar,\bbar]=[\pa(1)\pb(2)-\pb(1)\pa(2)]\beta_1\beta_2$ & $-1$\\
\bottomrule
\end{tabular}
\end{center}
Each entry follows from $D[\chi_1,\chi_2]=\chi_1(1)\chi_2(2)-\chi_2(1)\chi_1(2)$; for example
$D[a,\bbar]=\pa(1)\alpha_1\,\pb(2)\beta_2-\pb(1)\beta_1\,\pa(2)\alpha_2$.

\subsection{Triplet $M_S=\pm1$: single determinants}
Rows v and vi already factorise into (antisymmetric space)$\times$(symmetric spin $\Theta_{1,\pm1}$).
Normalising with $\tfrac1{\sqrt2}$ gives exactly the review's $^3\Psi(M_S=\pm1)$. Each is one determinant.
\textbf{[paper, confirmed]}

\subsection{$M_S=0$: two determinants are needed}
Determinants iii and iv individually are \emph{not} spin eigenfunctions: their spin parts
($\alpha_1\beta_2$ or $\beta_1\alpha_2$) are not of the forms \cref{eq:tripspin,eq:singspin}. Form the
sum and the difference. Write $A=\pa(1)\pb(2)$, $B=\pb(1)\pa(2)$:
\begin{align}
  D_{\text{iii}}&=A\,\alpha_1\beta_2-B\,\beta_1\alpha_2, \qquad
  D_{\text{iv}}=A\,\beta_1\alpha_2-B\,\alpha_1\beta_2.\\
  D_{\text{iii}}+D_{\text{iv}}&=A(\alpha_1\beta_2+\beta_1\alpha_2)-B(\beta_1\alpha_2+\alpha_1\beta_2)
     =(A-B)(\alpha_1\beta_2+\beta_1\alpha_2),\label{eq:sum}\\
  D_{\text{iv}}-D_{\text{iii}}&=A(\beta_1\alpha_2-\alpha_1\beta_2)-B(\alpha_1\beta_2-\beta_1\alpha_2)
     =(A+B)(\beta_1\alpha_2-\alpha_1\beta_2).\label{eq:diff}
\end{align}
Normalise: $\braket{A\mp B}{A\mp B}=2$ (orthonormal orbitals) and the spin factors have norm$^2=2$,
so each combination needs a factor $\tfrac12$.

\begin{resultbox}
\vspace{-1.2em}
\begin{align}
  ^3\Psi(M_S=0)&=\tfrac12\bigl(D_{\text{iii}}+D_{\text{iv}}\bigr)
     =\tfrac1{\sqrt2}[\pa(1)\pb(2)-\pb(1)\pa(2)]\;\tfrac1{\sqrt2}[\alpha_1\beta_2+\beta_1\alpha_2],\label{eq:T0}\\
  ^1\Psi_{\text{OSS}}&=\tfrac12\bigl(D_{\text{iv}}-D_{\text{iii}}\bigr)
     =\tfrac1{\sqrt2}[\pa(1)\pb(2)+\pb(1)\pa(2)]\;\tfrac1{\sqrt2}[\beta_1\alpha_2-\alpha_1\beta_2].\label{eq:OSS}
\end{align}
The open-shell singlet is irreducibly a two-determinant function; the high-spin triplet components are
single determinants. This is why a single-reference method describes the triplet but not the
open-shell singlet of a diradical. \textbf{[paper's physics, confirmed]}
\end{resultbox}

\begin{caveatbox}
The review writes the singlet as $\tfrac1{\sqrt2}\tfrac1{\sqrt2}(D_{\text{iii}}+D_{\text{iv}})$, with a
\emph{plus} sign. By \cref{eq:sum} that combination is the $M_S=0$ component of the \emph{triplet},
\cref{eq:T0}. The singlet requires the minus sign of \cref{eq:diff}. (Its overall sign is arbitrary;
$\tfrac12(D_{\text{iii}}-D_{\text{iv}})$ is equally correct.) The review's product form, in the line
above it, is correct. The physical conclusion is unaffected.
\end{caveatbox}

\paragraph{Closed-shell singlets.} Rows i and ii factorise as (symmetric space)$\times\sqrt2\,\Theta_{0,0}$:
\[
  \tfrac1{\sqrt2}D_{\text{i}}=\pa(1)\pa(2)\,\Theta_{0,0},\qquad \tfrac1{\sqrt2}D_{\text{ii}}=\pb(1)\pb(2)\,\Theta_{0,0}.
\]
So the two-electron/two-orbital ($2e/2o$) problem has \textbf{three singlets} (i, ii, OSS) and \textbf{one
triplet} (three components).

\begin{verifybox}
\texttt{t01\_spin\_and\_determinants.py}: every identity above expanded symbolically; the review's
``$+$'' combination is confirmed equal to $^3\Psi(M_S=0)$ and \emph{not} to $^1\Psi_{\text{OSS}}$. \textbf{PASS}.
\end{verifybox}

\section{Energies of the $2e/2o$ states and the full $2e/2o$ Hamiltonian}\label{sec:energies}

The review discusses Coulomb repulsion $J_{xy}$ and exchange $K_{xy}$ verbally (p.~442). Here are the
actual energy expressions.

\subsection{Matrix elements by direct integration}
For two electrons $\Hop=\hat h(1)+\hat h(2)+1/r_{12}$. Every state of \cref{sec:twoe} has the form
(spatial part $\Phi$)$\times$(normalised spin part $\Theta$). $\Hop$ is spin-free, so
$\mel{\Phi\Theta}{\Hop}{\Phi'\Theta'}=\braket{\Theta}{\Theta'}\,\mel{\Phi}{\Hop}{\Phi'}$. States with
different $S$ have orthogonal spin parts and do not mix. It is enough to evaluate spatial integrals.

\paragraph{Elementary integrals.} For orthonormal real orbitals:
\begin{align}
  \mel{\varphi_p(1)\varphi_q(2)}{\hat h(1)}{\varphi_r(1)\varphi_s(2)}&=h_{pr}\,\delta_{qs},\qquad
  \mel{\varphi_p(1)\varphi_q(2)}{\hat h(2)}{\varphi_r(1)\varphi_s(2)}=\delta_{pr}\,h_{qs},\label{eq:oneint}\\
  \mel{\varphi_p(1)\varphi_q(2)}{r_{12}^{-1}}{\varphi_r(1)\varphi_s(2)}&=\eri prqs .\label{eq:twoint}
\end{align}

\paragraph{Triplet.} $\Phi_T=\tfrac1{\sqrt2}(A-B)$ with $A=\pa(1)\pb(2)$, $B=\pb(1)\pa(2)$.
\begin{enumerate}
  \item One-electron part: $\mel{A}{\hat h(1)+\hat h(2)}{A}=h_{aa}+h_{bb}$, the same for $\mel{B}{\hat h(1)+\hat h(2)}{B}$, and
        $\mel{A}{\hat h(1)+\hat h(2)}{B}=h_{ab}\delta_{ba}+\delta_{ab}h_{ba}=0$. So the one-electron part
        is $\tfrac12(2h_{aa}+2h_{bb})=h_a+h_b$, where $h_a\equiv h_{aa}$.
  \item Two-electron part: $\mel A{r_{12}^{-1}}A=\eri aabb=J_{ab}$, $\mel B{r_{12}^{-1}}B=J_{ab}$, and
        $\mel A{r_{12}^{-1}}B=\eri abba=K_{ab}$. Then
        $\tfrac12\bigl(J_{ab}+J_{ab}-K_{ab}-K_{ab}\bigr)=J_{ab}-K_{ab}$.
\end{enumerate}
\begin{equation}\label{eq:ET}
  \boxed{E_T=h_a+h_b+J_{ab}-K_{ab}}
\end{equation}

\paragraph{Open-shell singlet.} $\Phi_{\text{OSS}}=\tfrac1{\sqrt2}(A+B)$. The same steps with $+$ signs
on the cross terms give
\begin{equation}\label{eq:EOSS}
  \boxed{E_{\text{OSS}}=h_a+h_b+J_{ab}+K_{ab}},\qquad E_{\text{OSS}}-E_T=2K_{ab}.
\end{equation}

\paragraph{Closed shells.} $\Phi_{\text i}=\pa(1)\pa(2)$: $E_{\text i}=2h_a+J_{aa}$, and
$E_{\text{ii}}=2h_b+J_{bb}$.

\paragraph{Couplings between the singlets.}
\begin{enumerate}
  \item $\mel{\Phi_{\text i}}{\Hop}{\Phi_{\text{ii}}}$: one-electron terms vanish ($\delta_{ab}=0$ factors);
        two-electron: $\mel{\pa\pa}{r_{12}^{-1}}{\pb\pb}=\eri abab=K_{ab}$.
  \item $\mel{\Phi_{\text i}}{\Hop}{\Phi_{\text{OSS}}}=\tfrac1{\sqrt2}\bigl[\mel{\pa\pa}{\Hop}{A}+\mel{\pa\pa}{\Hop}{B}\bigr]$.
        By \cref{eq:oneint},
        \[\mel{\pa(1)\pa(2)}{\hat h(1)+\hat h(2)}{\pa(1)\pb(2)}=0+h_{ab},\]
        and the same
        for $B$. By \cref{eq:twoint}, $\mel{\pa\pa}{r_{12}^{-1}}{A}=\eri aaab$ and
        $\mel{\pa\pa}{r_{12}^{-1}}{B}=\eri abaa=\eri aaab$. Sum: $\tfrac1{\sqrt2}\cdot2(h_{ab}+\eri aaab)$.
  \item Likewise $\mel{\Phi_{\text{ii}}}{\Hop}{\Phi_{\text{OSS}}}=\sqrt2\,(h_{ab}+\eri bbba)$.
\end{enumerate}

\begin{resultbox}
The singlet block of the exact $2e/2o$ Hamiltonian, in the basis $\{\Phi_{\text i},\Phi_{\text{ii}},\Phi_{\text{OSS}}\}$:
\begin{equation}\label{eq:Hsinglet}
\mathbf H_S=\begin{pmatrix}
2h_a+J_{aa} & K_{ab} & \sqrt2\,(h_{ab}+\eri aaab)\\
K_{ab} & 2h_b+J_{bb} & \sqrt2\,(h_{ab}+\eri bbba)\\
\sqrt2\,(h_{ab}+\eri aaab) & \sqrt2\,(h_{ab}+\eri bbba) & h_a+h_b+J_{ab}+K_{ab}
\end{pmatrix},
\end{equation}
and the triplet is the $1\times1$ block $E_T$ of \cref{eq:ET}. For orbitals of different symmetry
(e.g.\ $\sigma_g$, $\sigma_u$ of H$_2$; or the NBMOs of a symmetric diradical),
$h_{ab}=\eri aaab=\eri bbba=0$ and the OSS decouples from the closed shells.
\end{resultbox}

\paragraph{Consequence (Hund's rule at frozen orbitals).} With the \emph{same} orbitals in both states,
\cref{eq:EOSS} gives $E_{\text{OSS}}>E_T$ whenever $K_{ab}>0$. $K_{ab}\ge0$ is proved in
\cref{sec:holes}. This is Slater's argument: the triplet is lower by $2K_{ab}$ \emph{because} of a smaller
electron--electron repulsion ($J-K$ vs.\ $J+K$). \Cref{sec:hund} shows that this mechanism is
not the correct one once the orbitals are allowed to relax.

\begin{verifybox}
\texttt{t02\_two\_orbital\_hamiltonian.py}: for 5 random sets of integrals (with positive-definite
two-electron kernel), the eigenvalues of \cref{eq:Hsinglet} and $E_T$ equal the singlet and triplet roots
of a full-CI diagonalisation (PySCF \texttt{fci.direct\_spin1}, roots classified by $\ev{\hat S^2}$),
to $2\times10^{-14}$ (max.\ deviation over 5 sets); $\Phi_{\text{OSS}}$ diagonal minus $E_T$ equals $2K_{ab}$. \textbf{PASS}.
\end{verifybox}

\section{Fermi and Coulomb holes; positivity of $K$ and $J\ge K$}\label{sec:holes}

\begin{reviewclaim}
``\dots two electrons of the same spin cannot occupy the same space \dots\ The corresponding shortage
of the electron of a certain spin in the vicinity of another electron of the same spin is known as the
Fermi hole. Two opposite-spin electrons, on the other hand, are not subject to the exclusion principle.'' (p.~442)
\end{reviewclaim}

\subsection{Pair densities}
The spin-integrated pair density is $P(\mathbf r_1,\mathbf r_2)=N(N-1)\sum_{\sigma_1\sigma_2}\int|\Psi|^2\dd x_3\cdots$.
For two electrons it is $2\sum_{\sigma_1\sigma_2}|\Psi(x_1,x_2)|^2$. The spin parts are normalised, so
\begin{equation}
  P_T(\mathbf r_1,\mathbf r_2)=2\,|\Phi_T|^2=\bigl[\pa(1)\pb(2)-\pb(1)\pa(2)\bigr]^2,\qquad
  P_{\text{OSS}}=\bigl[\pa(1)\pb(2)+\pb(1)\pa(2)\bigr]^2 .
\end{equation}
\begin{enumerate}
  \item Set $\mathbf r_1=\mathbf r_2=\mathbf r$: $P_T(\mathbf r,\mathbf r)=[\pa\pb-\pb\pa]^2=0$ for \emph{every}
        $\mathbf r$. Two same-spin electrons are never found at the same point. This is the Fermi hole,
        and it is exact (no approximation beyond the form of $\Psi$).
  \item $P_{\text{OSS}}(\mathbf r,\mathbf r)=4\,\pa(\mathbf r)^2\pb(\mathbf r)^2\ge0$. It is not zero
        unless the orbitals do not overlap in space. There is no exclusion for an antiparallel pair,
        so the singlet's two electrons can come together. (The Coulomb hole, the reduction of $P$ at
        small $r_{12}$ caused by $1/r_{12}$ itself, is absent from these single-configuration functions.
        It appears only with correlation, i.e.\ configuration mixing.)
\end{enumerate}

\subsection{$K_{ab}\ge0$}
Let $\rho(\mathbf r)=\pa(\mathbf r)\pb(\mathbf r)$ (the overlap density). Then $K_{ab}=\iint\rho(1)\rho(2)/r_{12}$.
\begin{enumerate}
  \item Fourier representation of the Coulomb kernel: $\dfrac1{|\mathbf r|}=\displaystyle\int\frac{\dd^3k}{(2\pi)^3}\frac{4\pi}{k^2}\,e^{i\mathbf k\cdot\mathbf r}$.
  \item Insert this and define $\tilde\rho(\mathbf k)=\int\rho(\mathbf r)e^{-i\mathbf k\cdot\mathbf r}\dd\mathbf r$.
        Because $\rho$ is real, $\tilde\rho(-\mathbf k)=\tilde\rho(\mathbf k)^*$:
        \[
          K_{ab}=\int\frac{\dd^3k}{(2\pi)^3}\frac{4\pi}{k^2}\Bigl(\int\rho(1)e^{i\mathbf k\cdot\mathbf r_1}\Bigr)
          \Bigl(\int\rho(2)e^{-i\mathbf k\cdot\mathbf r_2}\Bigr)
          =\int\frac{\dd^3k}{(2\pi)^3}\frac{4\pi}{k^2}\,|\tilde\rho(\mathbf k)|^2 .
        \]
  \item The integrand is non-negative, so $K_{ab}\ge0$. Equality holds iff $\tilde\rho\equiv0$, i.e.\
        iff $\pa\pb=0$ everywhere (orbitals with no common region of space).
\end{enumerate}

\subsection{$J_{ab}\ge K_{ab}$}
\begin{enumerate}
  \item Take $f(1,2)=\pa(1)\pb(2)-\pb(1)\pa(2)$. Then $\iint f^2/r_{12}\ge0$ because the integrand is non-negative.
  \item Expand: $\iint f^2/r_{12}=J_{ab}+J_{ab}-2K_{ab}=2(J_{ab}-K_{ab})$ (same integrals as in \cref{sec:energies}).
  \item Hence $J_{ab}\ge K_{ab}\ge0$. The triplet repulsion $J-K$ is non-negative, and lies below the
        singlet value $J+K$.
\end{enumerate}
\begin{verifybox}
\texttt{t03\_t05\_hund\_virial\_h2.py} (T3): for all canonical RHF/6-31G orbitals of water,
$\min_{pq}K_{pq}=1.1\times10^{-3}\,E_h>0$ and $J_{pq}\ge K_{pq}$ for every pair. \textbf{PASS}.
\end{verifybox}

\section{Hund's multiplicity rule: Slater's argument, the virial theorem and ``less screening''}\label{sec:hund}

\begin{reviewclaim}
``Slater's original explanation of the Hund's rule preference for the high-spin state was therefore
based upon less electron--electron Coulomb repulsion in the triplet state \dots\ However, it has been
later shown that the electron--electron repulsion is actually greater in the high-spin state, and an
alternative explanation---the so-called less screening rule---has been offered. According to this rule,
Fermi holes in the triplet state lead to decreased screening of the nucleus and larger
electron--nucleus attraction. While this in turn causes the electron clouds to contract and partially
restores the electron--electron repulsion, the system overall is nonetheless stabilized.'' (p.~442,
refs.~11--14). Also: ``the absolute value of the correlation energy is approximately an order of
magnitude lower for a pair of two same-spin electrons compared to that of opposite-spin electrons'' (ref.~16).
\end{reviewclaim}

\subsection{Slater's argument (frozen orbitals)}
From \cref{eq:ET,eq:EOSS}, with identical orbitals in both states,
\begin{gather}\label{eq:slater}
  E_T-E_{\text{OSS}}=-2K_{ab}<0,\qquad
  \underbrace{\ev{V_{ee}}_T-\ev{V_{ee}}_{\text{OSS}}}_{=(J-K)-(J+K)}=-2K_{ab},\\
  \ev{T}_T-\ev{T}_{\text{OSS}}=0,\qquad \ev{V_{ne}}_T-\ev{V_{ne}}_{\text{OSS}}=0 .
\end{gather}
All of the stabilisation is assigned to $V_{ee}$.

\subsection{The virial theorem (scaling proof)}
\begin{proposition}\label{prop:virial}
For an exact bound eigenstate of an atom (one nucleus at the origin), $2\ev{T}+\ev{V}=0$, hence
$E=-\ev{T}=\tfrac12\ev V$. The same holds for any approximate wavefunction that is variationally optimal
with respect to uniform scaling.
\end{proposition}
\begin{proof}
\begin{enumerate}
  \item For a normalised $\Psi$ define $\Psi_\eta(\mathbf r_1,\dots,\mathbf r_N)=\eta^{3N/2}\Psi(\eta\mathbf r_1,\dots,\eta\mathbf r_N)$
        with $\eta>0$. Substituting $\mathbf s_i=\eta\mathbf r_i$ ($\dd^{3N}r=\eta^{-3N}\dd^{3N}s$) shows that
        $\Psi_\eta$ is normalised.
  \item Kinetic energy: $\nabla_{\mathbf r}=\eta\nabla_{\mathbf s}$, so $\ev{T}_\eta=\eta^2\ev T$.
  \item Potential energy: every Coulomb term ($-Z/r_i$ and $1/r_{ij}$) is homogeneous of degree $-1$,
        $1/r=\eta/s$, so $\ev V_\eta=\eta\ev V$.
  \item $E(\eta)=\eta^2\ev T+\eta\ev V$. The exact eigenstate is a stationary point of the energy
        functional with respect to \emph{all} variations, in particular this one, so
        $\left.\dd E/\dd\eta\right|_{\eta=1}=2\ev T+\ev V=0$.
  \item Then $E=\ev T+\ev V=\ev T-2\ev T=-\ev T$, and $E=\tfrac12\ev V$.
  \item (Molecules: the nuclear positions also scale. One finds $2\ev T+\ev V=-\sum_A\mathbf R_A\cdot\nabla_{\mathbf R_A}E$,
        which vanishes at an equilibrium geometry.)\qedhere
\end{enumerate}
\end{proof}

\subsection{Why the frozen-orbital argument cannot be the true mechanism}
\begin{enumerate}
  \item Apply \cref{prop:virial} to both exact states and subtract:
        \begin{equation}\label{eq:virialdiff}
          \Delta E\equiv E_T-E_S=-\Delta\ev T=\tfrac12\bigl(\Delta\ev{V_{ne}}+\Delta\ev{V_{ee}}\bigr).
        \end{equation}
  \item In \cref{eq:slater} $\Delta\ev T=0$. By \cref{eq:virialdiff} this would force $\Delta E=0$, contradicting
        $\Delta E=-2K<0$. The frozen-orbital pair of states violates the virial theorem, so it cannot
        describe how the \emph{exact} energy difference is partitioned. (This is precisely the objection of
        Hongo et al.: ``This assumption leaves the kinetic energy $T$ unchanged for different spin states,
        which cannot be correct as each eigenstate must satisfy the virial theorem.'')
  \item Since $\Delta E<0$, \cref{eq:virialdiff} demands $\Delta\ev T>0$ (the triplet electrons are more
        compressed) and $\Delta\ev V=2\Delta E<0$. If, as the review states, $\Delta\ev{V_{ee}}>0$, then necessarily
        \begin{equation}\label{eq:lessscreen}
          \Delta\ev{V_{ne}}=2\Delta E-\Delta\ev{V_{ee}}<2\Delta E<0,\qquad|\Delta\ev{V_{ne}}|>\Delta\ev{V_{ee}} .
        \end{equation}
        The triplet is stabilised by a stronger electron--nucleus attraction that over-compensates the
        larger electron--electron repulsion: the less-screening mechanism.
  \item Physical picture: the Fermi hole (\cref{sec:holes}) keeps parallel electrons apart, so each
        electron screens the nucleus less from the other. The electrons contract towards the nucleus,
        which lowers $V_{ne}$, raises $T$ (\cref{eq:virialdiff}), and also raises $V_{ee}$, because
        contracted electrons are on average closer together.
\end{enumerate}

\subsection{Numerical test: helium $1s2s$, $2\,^3S$ versus $2\,^1S$}
Helium is the simplest case: two electrons, and the same configuration $1s2s$ gives a singlet and a triplet.
The two-electron Schr\"odinger equation was solved exactly within a large even-tempered Gaussian basis
($16s7p4d$). The spatial function $\Psi(1,2)=\sum_{pq}C_{pq}\varphi_p(1)\varphi_q(2)$ is restricted to symmetric
$C$ (singlet) or antisymmetric $C$ (triplet), following \cref{sec:anti}. The exponents were uniformly scaled
by the variationally optimal $\eta^2$, which by \cref{prop:virial} enforces the virial theorem.
\begin{center}
\begin{tabular}{lrrrrr}
\toprule
state & $E$ & exact $E$ (Drake) & $\ev T$ & $\ev{V_{ne}}$ & $\ev{V_{ee}}$\\
\midrule
$2\,^3S$ & $-2.175212$ & $-2.175229$ & $2.175212$ & $-4.618633$ & $0.268208$\\
$2\,^1S$ & $-2.145865$ & $-2.145974$ & $2.145865$ & $-4.541465$ & $0.249735$\\
\midrule
$\Delta(T-S)$ & $-0.029347$ & $-0.029255$ & $+0.029347$ & $-0.077168$ & $+0.018473$\\
\bottomrule
\end{tabular}
\end{center}
The virial ratio $-\ev V/\ev T=2.000000$ for both states. Check of \cref{eq:virialdiff}:
$\tfrac12(-0.077168+0.018473)=-0.029348\approx\Delta E$. The triplet has the \emph{larger} electron
repulsion ($+0.0185\,E_h$), the opposite of Slater's picture. It is stabilised by a $0.077\,E_h$ gain in
nuclear attraction, exactly as in \cref{eq:lessscreen}.

\paragraph{Literature values for carbon ($^3P$ vs.\ $^1D$), Hongo et al.\ 2004, Table I [verified: arXiv cond-mat/0408147].}
\begin{center}
\begin{tabular}{lrrr}
\toprule
method & $\Delta\ev{V_{ee}}$ & $\Delta\ev{V_{en}}$ & $\Delta\ev T$\\
\midrule
HF & $+0.032$ & $-0.145$ & $+0.057$\\
DMC (linear estimator) & $+0.074$ & $-0.169$ & $+0.048$\\
\bottomrule
\end{tabular}
\end{center}
(Values in $E_h$, triplet minus singlet, computed from their Table~I.) The pattern is the same as for He.

\subsection{Same-spin versus opposite-spin correlation}
The review (citing Krylov 2001) states that the correlation energy of a same-spin pair is about an order
of magnitude smaller than that of an opposite-spin pair. Krylov's introduction says this
[verified: ScienceDirect preview of Chem.\ Phys.\ Lett.\ 338, 375]. The reason is the Fermi hole: the
same-spin pair is already kept apart at the mean-field level, so there is less left for correlation to do.
Tests:
\begin{itemize}
  \item He: $E_{\text{corr}}(1\,^1S)=-0.04205$ (one $\uparrow\downarrow$ pair) vs.\ $E_{\text{corr}}(2\,^3S)=-0.00098\,E_h$
        (one $\uparrow\uparrow$ pair; exact minus ROHF in the same basis). Ratio $0.023$: the statement holds,
        and strongly.
  \item Molecules, MP2/cc-pVTZ, frozen core: same-spin/opposite-spin ratio of the total correlation energy
        is $0.32$ (H$_2$O), $0.35$ (N$_2$), $0.27$ (C$_2$H$_4$). \emph{Per pair} it is $0.43$, $0.43$, $0.32$.
\end{itemize}
\begin{caveatbox}
``An order of magnitude'' is right for compact two-electron systems such as He, but overstated for
valence electrons in molecules at the MP2 level, where same-spin pairs carry roughly one third of the
correlation per pair. This is consistent with the empirical SCS-MP2 scaling ($c_{os}=6/5$, $c_{ss}=1/3$; Grimme 2003).
\end{caveatbox}

\begin{verifybox}
T4 (\texttt{t03\_t05\_hund\_virial\_h2.py}): energies within $1.7\times10^{-5}$ and $1.1\times10^{-4}\,E_h$ of the exact
values; virial satisfied; $\Delta V_{ee}>0$, $\Delta V_{ne}<0$, $|\Delta V_{ne}|>\Delta V_{ee}$. \textbf{PASS}.
T15 (\texttt{t09b\_t15\_correlation.py}): numbers above. \textbf{PASS}.
\end{verifybox}
\begin{sourcebox}
Boyd, \emph{Nature} \textbf{310}, 480 (1984), doi:10.1038/310480a0. Abstract read: Hund's rule is a
consequence of Fermi correlation, which leads to less screening of the nucleus in the high-spin state.
Hongo et al., \emph{J.\ Chem.\ Phys.}\ \textbf{121}, 7144 (2004): full text read (arXiv). Katriel \& Pauncz,
\emph{Adv.\ Quantum Chem.}\ \textbf{10}, 143 (\textbf{1977}, not 1975 as cited), doi:10.1016/S0065-3276(08)60580-8:
not accessed. Hongo et al., \emph{J.\ Magn.\ Magn.\ Mater.}\ \textbf{310}, e560 (2007): not accessed.
Slater (1929): cited by the review only via Kollmar \& Staemmler (ref.~9).
\end{sourcebox}

\section{Static correlation: H$_2$ dissociation and the weak-interaction limit (Figure 2)}\label{sec:static}

\begin{reviewclaim}
``\dots dynamic correlation is weak compared to static correlation effects, i.e.\ correlation between
configurations in the multireference wavefunction \dots\ The weaker the interaction between the two
nonbonding molecular orbitals, the stronger the static correlation (Figure 2). This is clearly
illustrated by \dots\ the H$_2$ homolytic dissociation curve or the singlet ground state of square
cyclobutadiene \dots'' (p.~442)
\end{reviewclaim}

\subsection{H$_2$ in a minimal basis: exact two-configuration treatment}
Basis: one $1s$ function on each nucleus, $\chi_A,\chi_B$, overlap $S$. By symmetry the MOs are
\begin{equation}
  \sigma_g\equiv\varphi_1=\frac{\chi_A+\chi_B}{\sqrt{2(1+S)}},\qquad
  \sigma_u\equiv\varphi_2=\frac{\chi_A-\chi_B}{\sqrt{2(1-S)}} .
\end{equation}
$\varphi_1$ is even and $\varphi_2$ odd under inversion, so $h_{12}=\eri1112=\eri2221=0$ and \cref{eq:Hsinglet}
splits. The OSS ($^1\Sigma_u^+$) is alone, and the two closed shells ($^1\Sigma_g^+$ symmetry) mix:
\begin{equation}\label{eq:h2ci}
  \mathbf H(^1\Sigma_g^+)=\begin{pmatrix}2h_{11}+J_{11}&K_{12}\\K_{12}&2h_{22}+J_{22}\end{pmatrix},\qquad
  E_{\text{CI}}=\tfrac12\Bigl[E_1+E_2-\sqrt{(E_2-E_1)^2+4K_{12}^2}\Bigr],
\end{equation}
with $E_1=2h_{11}+J_{11}$ (the RHF energy), $E_2=2h_{22}+J_{22}$, and
$E(^3\Sigma_u^+)=h_{11}+h_{22}+J_{12}-K_{12}$, $E(^1\Sigma_u^+)=h_{11}+h_{22}+J_{12}+K_{12}$
(nuclear repulsion $1/R$ added to all).

\subsection{The dissociation limit, step by step}
Let $R\to\infty$, so $S\to0$ and all two-centre overlap densities $\chi_A\chi_B\to0$. Define the on-site
repulsion $U=\eri AAAA$, the inter-site repulsion $V=\eri AABB\to1/R$, and the H-atom energy $E_H$.
\begin{enumerate}
  \item Orbital products: $\varphi_1^2=\tfrac12(\chi_A^2+\chi_B^2+2\chi_A\chi_B)\to\tfrac12(\chi_A^2+\chi_B^2)$,
        $\varphi_2^2\to\tfrac12(\chi_A^2+\chi_B^2)$, $\varphi_1\varphi_2\to\tfrac12(\chi_A^2-\chi_B^2)$.
  \item Two-electron integrals (each product contributes $\tfrac12\cdot\tfrac12$):
        \begin{align*}
        J_{11}=J_{22}=J_{12}&\to\tfrac14\bigl[\eri AAAA+\eri AABB+\eri BBAA+\eri BBBB\bigr]=\tfrac12(U+V),\\
        K_{12}=\eri1212&\to\tfrac14\bigl[\eri AAAA-\eri AABB-\eri BBAA+\eri BBBB\bigr]=\tfrac12(U-V).
        \end{align*}
  \item One-electron integrals: $h_{11}=h_{22}\to h_{AA}=E_H-\tfrac1R$ (an electron on A also feels nucleus B
        as a point charge at distance $R$).
  \item RHF: $E_{\text{RHF}}=2h_{11}+J_{11}+\tfrac1R=2E_H-\tfrac2R+\tfrac12(U+V)+\tfrac1R\to2E_H+\tfrac U2-\tfrac1{2R}$.
        \emph{The RHF error at dissociation is $U/2$}, which does not vanish.
  \item Triplet: $E_T=2h_{11}+J_{12}-K_{12}+\tfrac1R=2E_H-\tfrac2R+V+\tfrac1R\to2E_H$. Correct.
  \item CI: the diagonal elements are now equal ($E_1=E_2$), so by \cref{eq:h2ci}
        $E_{\text{CI}}=E_1-K_{12}=2E_H+\tfrac U2-\tfrac1{2R}-\tfrac12(U-V)\to2E_H$. Correct.
        The eigenvector has $c_1=-c_2=\tfrac1{\sqrt2}$: \emph{two configurations with equal weight}. This is
        maximal static correlation.
  \item What that eigenvector is: $\varphi_1(1)\varphi_1(2)-\varphi_2(1)\varphi_2(2)\to\tfrac12[(\chi_A+\chi_B)(\chi_A+\chi_B)-(\chi_A-\chi_B)(\chi_A-\chi_B)]
        =\chi_A(1)\chi_B(2)+\chi_B(1)\chi_A(2)$, the purely covalent Heitler--London function. By contrast
        $\varphi_1(1)\varphi_1(2)$ alone is $\tfrac12[\text{covalent}+\text{ionic}]$, which is the origin of the
        $U/2$ error in step 4.
  \item OSS: $E(^1\Sigma_u^+)=E_T+2K_{12}\to2E_H+U-V$. It is the antisymmetric \emph{ionic} combination.
\end{enumerate}

\begin{verifybox}
T5 (\texttt{t03\_t05\_hund\_virial\_h2.py}), STO-3G, PySCF vs.\ the closed forms at $R=1.4,3,6,20$ bohr.
RHF, 2$\times$2 CI\,=\,FCI and ROHF triplet agree with \cref{eq:h2ci} to $4\times10^{-16}\,E_h$ (max.\ deviation, tight SCF). At $R=20$:
$E_{\text{CI}}=E_T=-0.933164=2E_H$; $E_{\text{RHF}}-2E_H=0.3623=K_{12}$; $J_{11}-K_{12}=0.050=1/R$ (so $U=J_{11}+K_{12}=0.7746$);
$E(^1\Sigma_u^+)-2E_H=0.7246=U-1/R$. Every limit in steps 4--8 is reproduced. \textbf{PASS}.
\end{verifybox}

\subsection{General weak-interaction limit (Figure 2b, twisted ethylene)}
Take the $2e/2o$ model with orbitals of different symmetry ($h_{ab}=\eri aaab=\eri bbba=0$), orbital
energy gap $\Delta$ ($h_a=-\Delta/2$, $h_b=+\Delta/2$), and for simplicity $J_{aa}=J_{bb}\equiv J_d$,
$J_{ab}\equiv J_x$. This describes the $\pi/\pi^*$ pair of ethylene, where $\Delta$ is the $\pi$--$\pi^*$
splitting, which goes to $0$ at a $90^\circ$ twist.
\begin{enumerate}
  \item Closed-shell block: $\begin{pmatrix}J_d-\Delta&K\\K&J_d+\Delta\end{pmatrix}$, eigenvalues $J_d\pm\sqrt{\Delta^2+K^2}$.
  \item Ground singlet: $E_S=J_d-\sqrt{\Delta^2+K^2}$, with coefficient ratio
        $\dfrac{c_2}{c_1}=\dfrac{E_S-(J_d-\Delta)}{K}=\dfrac{\Delta-\sqrt{\Delta^2+K^2}}{K}$.
  \item Limits: $\Delta\gg K$ gives $c_2/c_1\to0$ (one configuration, closed shell, no static correlation);
        $\Delta\to0$ gives $c_2/c_1\to-1$ (two configurations of equal weight, a pure diradical).
  \item Gap to the triplet ($E_T=J_x-K$): $E_S-E_T=(J_d-J_x)+K-\sqrt{\Delta^2+K^2}$.
\end{enumerate}
\begin{resultbox}
Static correlation (the weight of the second configuration) is controlled by $\Delta/K$. The smaller
the splitting of the two frontier orbitals (``weaker interaction''), the more multiconfigurational the
singlet becomes, and the closer it comes to the triplet. This is the content of Figure~2(a,b).
\end{resultbox}
\begin{verifybox}
T6 (\texttt{t06\_t09\_models.py}): eigenvalue, ratio and both limits verified symbolically. \textbf{PASS}.
Balkov\'a \& Bartlett, \emph{J.\ Chem.\ Phys.}\ \textbf{101}, 8972 (1994), doi:10.1063/1.468057
[abstract]: MR-CC places the square-cyclobutadiene singlet 6.9~kcal/mol below the triplet.
\end{verifybox}
Figure~2(c) (spin crossover in a transition-metal complex) is qualitative in the review and contains no
mathematics to derive.

\section{Localised orbitals, ionic configurations and kinetic exchange}\label{sec:kinex}

\begin{reviewclaim}
``The kinetic exchange picture based on the Hubbard model is a related, yet also simplistic (and not
satisfying virial theorem) spin alignment rule. It attributes stability of a singlet state in the case
of orthogonal atomic orbitals (AOs) to the fact that perturbation yielding ionic configuration is
allowed only from this state, and not from the triplet.'' (p.~442, ref.~19 Yamanaka et al.\ 2005).
Also: ``the importance of resonance between open-shell singlet and closed-shell zwitterionic
configurations'' (p.~444, Figure~6).
\end{reviewclaim}

\subsection{Two-site model in a localised basis}
Two orthonormal localised orbitals $l,r$ (one per radical centre). Parameters:
\begin{gather}
  h_{ll}=h_{rr}=0,\quad h_{lr}=-t,\quad U=\eri llll=\eri rrrr,\quad V=\eri llrr,\\
  k=\eri lrlr\ (\ge0),\quad \eri lllr=\eri rrrl=0 .
\end{gather}
($V=k=0$ is the Hubbard dimer; $V,k\neq0$ the ``extended'' model.)

\paragraph{Configurations.} \Cref{sec:twoe} with $a\to l$, $b\to r$: covalent singlet
$\Phi_{\text{cov}}=\tfrac1{\sqrt2}[l(1)r(2)+r(1)l(2)]\Theta_{0,0}$; ionic $l(1)l(2)\Theta_{0,0}$ and $r(1)r(2)\Theta_{0,0}$;
triplet $\tfrac1{\sqrt2}[l(1)r(2)-r(1)l(2)]\Theta_{1,M}$.

\paragraph{Step 1: the triplet cannot mix with ionic configurations.} An ionic configuration has both
electrons in the same spatial orbital, so its spatial part $l(1)l(2)$ is symmetric. By \cref{sec:anti} it
must carry the singlet spin function. Since $\Hop$ is spin-free, $\mel{\text{ionic}}{\Hop}{\text{triplet}}=0$
(orthogonal spin parts). This is the precise meaning of ``perturbation yielding ionic configuration is
allowed only from'' the singlet.

\paragraph{Step 2: matrix elements.} Insert the parameters into \cref{eq:Hsinglet,eq:ET}:
\begin{equation}
  E_T=V-k,\qquad
  \mathbf H_S=\begin{pmatrix}U&k&-\sqrt2t\\k&U&-\sqrt2t\\-\sqrt2t&-\sqrt2t&V+k\end{pmatrix}
  \ \ (\text{basis } ll,\ rr,\ \text{cov}).
\end{equation}

\paragraph{Step 3: symmetry-adapt the ionic pair.} $\Phi_{\text{ion}\pm}=\tfrac1{\sqrt2}(ll\pm rr)$ has
energy $U\pm k$. Its coupling to cov is $\tfrac1{\sqrt2}(-\sqrt2t\mp\sqrt2t)$, i.e.\ $-2t$ for ion$+$ and
$0$ for ion$-$. The ground singlet comes from the $2\times2$ block
$\begin{pmatrix}V+k&-2t\\-2t&U+k\end{pmatrix}$:
\begin{equation}\label{eq:ES_ext}
  E_S=\tfrac12\Bigl[U+V+2k-\sqrt{(U-V)^2+16t^2}\Bigr].
\end{equation}

\paragraph{Step 4: singlet--triplet gap and its expansion.}
\begin{equation}\label{eq:gap_ext}
  E_S-E_T=\tfrac12(U-V)+2k-\tfrac12\sqrt{(U-V)^2+16t^2}
  \;=\;2k-\frac{4t^2}{U-V}+O\!\left(\frac{t^4}{(U-V)^3}\right),
\end{equation}
using $\sqrt{x^2+16t^2}=x+8t^2/x+O(t^4/x^3)$ for $x=U-V\gg t$.

\begin{resultbox}
Two competing terms:
\begin{itemize}
  \item $+2k$, \emph{potential (direct) exchange}. It favours the triplet (Hund-like, \cref{sec:energies}).
  \item $-4t^2/(U-V)$, \emph{kinetic exchange}. The covalent singlet lowers its energy by virtual hopping into the
        ionic configuration (second-order perturbation: $|{-2t}|^2/[(V+k)-(U+k)]$). The triplet cannot do this
        (step~1), so it favours the singlet.
\end{itemize}
Mapped onto $\hat H_{\text{spin}}=-2J\,\hat{\mathbf S}_1\cdot\hat{\mathbf S}_2$ (for which $E_S-E_T=2J$):
$J\simeq k-2t^2/(U-V)$. For the pure Hubbard dimer ($V=k=0$), $E_S=\tfrac12(U-\sqrt{U^2+16t^2})\simeq-4t^2/U$
and $E_T=0$, i.e.\ antiferromagnetic coupling $J_{\text{AF}}=4t^2/U$ in the convention
$\hat H=J_{\text{AF}}\hat{\mathbf S}_1\cdot\hat{\mathbf S}_2$.
\end{resultbox}

\paragraph{Ionic resonance in charged polyradicals (Figure 6).} Pyridinium-bridged diradicals (V in
Figure~6) are the case where the closed-shell zwitterionic (``ionic'') resonance structures are low in
energy: effectively $U-V$ is small. By \cref{eq:gap_ext} the kinetic term $-4t^2/(U-V)$ then grows and
can overcome $2k$, reversing the topology-predicted triplet. This is the review's point that this
contribution ``exceeds the exchange and Coulomb terms''. (A qualitative mapping; the review gives no
numbers.)

\paragraph{On ``not satisfying the virial theorem''.} The Hubbard model has no separate $\hat T$ and $\hat V$.
Its energy is not a stationary point with respect to coordinate scaling, so \cref{prop:virial} cannot be
imposed. Like the frozen-orbital model of \cref{sec:hund}, it cannot say which physical energy component
pays for the spin preference.

\paragraph{Relation to delocalised orbitals.} With $a,b=(l\pm r)/\sqrt2$ one finds
$a(1)a(2)-b(1)b(2)=l(1)r(2)+r(1)l(2)$ (as in \cref{sec:static}, step 7). The ``two-configuration'' singlet of
the delocalised picture is the covalent singlet of the localised picture. Static correlation in one
basis is the absence of ionic character in the other.

\begin{verifybox}
T7 (\texttt{t06\_t09\_models.py}): \cref{eq:ES_ext,eq:gap_ext} verified symbolically, including the series
coefficient $-4t^2/U$; and numerically against PySCF full CI with the same integrals ($t=1$, $U=6$,
$V=1.5$, $k=0.2$), and against an independent fermionic exact-diagonalisation code. \textbf{PASS}.
\end{verifybox}
\begin{sourcebox}
Yamanaka et al., \emph{Int.\ J.\ Quantum Chem.}\ \textbf{105}, 687 (2005): not accessed. The derivation above
is the standard Anderson kinetic-exchange argument, reconstructed independently and tested.
\end{sourcebox}

\section{Non-bonding MOs, disjointness and the sign of $K$ (Borden--Davidson)}\label{sec:huckel}

\begin{reviewclaim}
``Borden and Davidson proposed to classify the nonbonding molecular orbitals (NBMOs) $\varphi_a$ and
$\varphi_b$ as either disjoint or nondisjoint. If the two orbitals do not have common atoms and the
probability of two electrons occupying the same space is nil independent of their spins, they are
termed disjoint. Under such circumstances, the energy of an open-shell singlet can become comparable to
or even lower than that of a triplet.'' (p.~443) Examples: CBD, TME (disjoint, singlet); MX, TriMB
(nondisjoint, high spin) (p.~443--444, Figure~3).
\end{reviewclaim}

\subsection{Hückel NBMOs: the zero-sum rule}
Hückel Hamiltonian for a conjugated $\pi$ system with adjacency matrix $\mathbf A$ ($A_{ij}=1$ if atoms
$i,j$ are bonded): $\mathbf H=\alpha\mathbf 1+\beta\mathbf A$. Writing $E=\alpha+x\beta$, the secular equation
is $\mathbf A\mathbf c=x\mathbf c$.
\begin{enumerate}
  \item An NBMO has $E=\alpha$, i.e.\ $x=0$, so $\mathbf A\mathbf c=\mathbf 0$.
  \item Row $i$ of this equation reads $\sum_{j\in\mathcal N(i)}c_j=0$. \textbf{Zero-sum rule:} for every atom, the
        NBMO coefficients of its neighbours sum to zero.
\end{enumerate}

\subsection{Alternant hydrocarbons: pairing theorem and NBMO count}
If the atoms split into ``starred'' ($\mathcal A$, $n_A$ atoms) and ``unstarred'' ($\mathcal B$, $n_B$ atoms) sets with
bonds only between the sets (a bipartite graph), then in that ordering
$\mathbf A=\begin{pmatrix}\mathbf0&\mathbf B\\\mathbf B^{\mathsf T}&\mathbf0\end{pmatrix}$, with $\mathbf B$ of size $n_A\times n_B$.
\begin{enumerate}
  \item \emph{Pairing theorem.} If $\mathbf A\binom{\mathbf u}{\mathbf v}=x\binom{\mathbf u}{\mathbf v}$, i.e.\
        $\mathbf B\mathbf v=x\mathbf u$ and $\mathbf B^{\mathsf T}\mathbf u=x\mathbf v$, then
        $\mathbf A\binom{\mathbf u}{-\mathbf v}=\binom{-\mathbf B\mathbf v}{\mathbf B^{\mathsf T}\mathbf u}=\binom{-x\mathbf u}{x\mathbf v}=-x\binom{\mathbf u}{-\mathbf v}$.
        Levels come in $\pm x$ pairs, with coefficients equal in magnitude and the sign flipped on one set.
  \item \emph{Lower bound from the sublattice imbalance.} $\operatorname{rank}\mathbf A=2\operatorname{rank}\mathbf B\le2\min(n_A,n_B)$.
        Hence the number of NBMOs (the nullity) is $N-\operatorname{rank}\mathbf A\ge N-2\min(n_A,n_B)=|n_A-n_B|$.
  \item \emph{Longuet-Higgins bound.} $\operatorname{rank}\mathbf B$ cannot exceed the largest number of
        non-attacking nonzero entries of $\mathbf B$. That is the size $T$ of a maximum matching, which equals
        the number of double bonds in the best Kekulé structure. So $\#\text{NBMO}\ge N-2T\ \ (\ge|n_A-n_B|)$.
        \textbf{This is a lower bound, not an equality}: square cyclobutadiene has $T=2$, so $N-2T=0$,
        yet it has two NBMOs (its $\pm x$ pairs include an accidental $x=0$ pair).
\end{enumerate}

\subsection{Which NBMO pairs are disjoint? Explicit constructions}
\paragraph{TME (atoms 1,2 central; 3,4 on C1; 5,6 on C2).} Zero-sum at atoms 3 and 4 (single neighbour 1)
gives $c_1=0$. At atoms 5 and 6 it gives $c_2=0$. At atom 1: $c_2+c_3+c_4=0$, so $c_4=-c_3$. At atom 2:
$c_5=-c_6$. The NBMO space is spanned by
\begin{equation}
  \varphi_a=\tfrac1{\sqrt2}(\chi_3-\chi_4),\qquad\varphi_b=\tfrac1{\sqrt2}(\chi_5-\chi_6),
\end{equation}
which share no atom: \textbf{disjoint} (two allyl-type NBMOs on separate halves).

\paragraph{Square CBD (ring 1-2-3-4).} Zero-sum at atom 1: $c_2+c_4=0$; at atom 2: $c_1+c_3=0$. So
$\varphi_a=\tfrac1{\sqrt2}(\chi_1-\chi_3)$ and $\varphi_b=\tfrac1{\sqrt2}(\chi_2-\chi_4)$: \textbf{disjoint}.

\paragraph{TMM (centre 0, terminals 1,2,3).} $c_0=0$ (zero-sum at a terminal), and $c_1+c_2+c_3=0$ (zero-sum
at the centre). Every nonzero vector in this plane has at least two nonzero entries (one nonzero entry
cannot sum to zero). Two such vectors with disjoint supports would need at least $2+2=4$ terminal
atoms, but there are only 3. So no basis of the NBMO space is disjoint: \textbf{nondisjoint}.

\subsection{Why disjointness kills $K_{ab}$: zero differential overlap}
Expand $\varphi_a=\sum_\mu a_\mu\chi_\mu$, $\varphi_b=\sum_\nu b_\nu\chi_\nu$ in atomic $p_\pi$ orbitals:
\begin{enumerate}
  \item $K_{ab}=\eri abab=\sum_{\mu\nu\lambda\sigma}a_\mu b_\nu a_\lambda b_\sigma\eri{\mu}{\nu}{\lambda}{\sigma}$.
  \item ZDO approximation: $\eri\mu\nu\lambda\sigma=\delta_{\mu\nu}\delta_{\lambda\sigma}\gamma_{\mu\lambda}$, with $\gamma_{\mu\lambda}$ the
        Coulomb repulsion between charges on atoms $\mu$ and $\lambda$.
  \item Then $K_{ab}=\sum_{\mu\lambda}(a_\mu b_\mu)(a_\lambda b_\lambda)\gamma_{\mu\lambda}$.
  \item \emph{Disjoint} $\Rightarrow$ $a_\mu b_\mu=0$ for every atom $\Rightarrow$ $K_{ab}=0$ exactly (within ZDO).
        By \cref{eq:EOSS}, $E_{\text{OSS}}-E_T=2K_{ab}=0$: the first-order Hund preference vanishes, and
        higher-order effects (spin polarisation, \cref{sec:spinpol}) decide the ground state.
  \item \emph{Nondisjoint} $\Rightarrow$ some $a_\mu b_\mu\neq0$. Since $\gamma$ is a positive-definite Coulomb
        matrix (\cref{sec:holes}), $K_{ab}>0$ and the triplet is favoured at first order.
  \item In the Hubbard limit ($\gamma_{\mu\lambda}=U\delta_{\mu\lambda}$): $K_{ab}=U\sum_\mu a_\mu^2b_\mu^2$.
\end{enumerate}

\begin{verifybox}
T8 (\texttt{t06\_t09\_models.py}): Hückel spectra of CBD, TMM, TME, MX, TriMB, 1,2,4,5-TetraMB and
$p$-xylylene. Pairing theorem, zero-sum rule and $\#\text{NBMO}\ge N-2T\ge|n_A-n_B|$ hold in every case;
CBD shows the strict inequality. Minimising $\sum_\mu a_\mu^2b_\mu^2$ over all rotations of the degenerate
NBMO pair gives $\approx0$ for CBD, TME and TetraMB (disjoint), and $1/6$ (TMM) and $1/16$ (MX) (nondisjoint).
T9b (\texttt{t09b\_t15\_correlation.py}): in the Hubbard model at small $U$ the triplet--singlet gap of TMM
and MX is linear in $U$ (fitted exponent 0.97, 0.96). For TMM $-0.01643$ at $U=0.05t$ matches
$-2U\cdot\tfrac16=-0.0167$, and MX matches $-2U\cdot\tfrac1{16}$. For disjoint CBD and TME the gap is
$\propto U^{1.97},U^{1.95}$, i.e.\ zero at first order, as derived. \textbf{PASS}.
\end{verifybox}
\begin{sourcebox}
Borden \& Davidson, \emph{J.\ Am.\ Chem.\ Soc.}\ \textbf{99}, 4587 (1977), doi:10.1021/ja00456a010: not accessed;
definition as quoted in the review. Longuet-Higgins, \emph{J.\ Chem.\ Phys.}\ \textbf{18}, 265 (1950): the matching
bound as cited in arXiv:1906.08544. The proofs above are self-contained.
\end{sourcebox}

\section{Dynamic spin polarisation, topology rules, Ovchinnikov and Lieb}\label{sec:spinpol}

\begin{reviewclaim}
``\dots instantaneous correlation effects, in particular the dynamic spin polarization (Scheme 1). A
damped spin wave so formed generates ferromagnetic coupling between the spins separated by an odd number
of sites, and antiferromagnetic between those separated by an even number of sites, thus yielding the
topology rules of spin coupling. Within the spin-Hamiltonian VB theory, Ovchinnikov's qualitative rule
dictates that the most spin-alternant determinant is expected to have the lowest energy and the highest
contribution.'' (p.~443)
\end{reviewclaim}

\subsection{Spin polarisation is a correlation effect}
\begin{enumerate}
  \item In Hückel (independent-electron) theory the unpaired electron of an odd alternant radical sits in
        the NBMO. By \cref{sec:huckel} the NBMO vanishes on the smaller (unstarred) set. The spin density on
        unstarred atoms is therefore \emph{exactly zero}.
  \item Add on-site repulsion $U\sum_i\hat n_{i\uparrow}\hat n_{i\downarrow}$ (Hubbard model). An $\uparrow$ electron
        in the NBMO repels $\downarrow$ density on the starred atoms more than $\uparrow$ density. Paired
        electrons therefore rearrange so that $\downarrow$ density accumulates on the neighbouring (unstarred)
        atoms: a negative spin density appears between positive ones. This is Scheme~1 of the review.
\end{enumerate}
\begin{verifybox}
T9: pentadienyl radical (5-site chain, 5 electrons, $S_z=\tfrac12$), exact diagonalisation.
$U=0$: $\rho_s=(0.333,0,0.333,0,0.333)$. $U=3t$: $\rho_s=(0.451,-0.182,0.462,-0.182,0.451)$, sum $=1$.
The alternation appears only with $U\neq0$. \textbf{PASS}.
\end{verifybox}

\subsection{Odd/even rule from the bipartite structure}
\begin{enumerate}
  \item From \cref{sec:kinex}, for $U\gg t$ each bond carries an antiferromagnetic coupling
        $J_{\text{AF}}=4t^2/U$. At half filling the low-energy physics is the Heisenberg model
        $\hat H=J_{\text{AF}}\sum_{\langle ij\rangle}\hat{\mathbf S}_i\cdot\hat{\mathbf S}_j$.
  \item On a bipartite graph every bond joins $\mathcal A$ to $\mathcal B$. The energy is lowest when neighbours
        are antiparallel: spins on $\mathcal A$ point one way and spins on $\mathcal B$ the other (the Néel pattern).
  \item Two sites are in the same set iff the path between them has an even number of bonds, i.e.\ an
        \emph{odd} number of intervening atoms. Hence:
        same set $\Rightarrow$ parallel (ferromagnetic); different sets $\Rightarrow$ antiparallel.
  \item Example: in $m$-xylylene the radical CH$_2$ carbons are joined via 3 intervening ring atoms (odd), so
        they couple ferromagnetically (triplet). In $p$-xylylene there are 4 (even), so the coupling is
        antiferromagnetic: a singlet (actually Kekulé, closed shell). This is the review's rule.
\end{enumerate}

\subsection{Ovchinnikov's rule and Lieb's theorem}
\begin{enumerate}
  \item Counting the Néel configuration: $n_A$ spins up and $n_B$ spins down gives $S=\tfrac12|n_A-n_B|$.
        This is \textbf{Ovchinnikov's rule} for the ground-state spin of an alternant hydrocarbon,
        $S=\tfrac12|n^*-n^\circ|$.
  \item The rule was conjectured from VB/Heisenberg arguments. Lieb (1989) proved it rigorously for the
        repulsive Hubbard model at half filling on any bipartite graph: \emph{the ground state is unique and
        has $S=\tfrac12\bigl||\mathcal B|-|\mathcal A|\bigr|$}, for every $U>0$ [verified: \emph{Phys.\ Rev.\ Lett.}\ \textbf{62}, 1201 (1989)].
  \item Consequences: CBD and TME ($n_A=n_B$) have singlet ground states; TMM (1/3) and MX (3/5) are triplets;
        TriMB (3/6) is a quartet; 1,2,4,5-tetramethylenebenzene (5/5) is a singlet. These are the review's
        Figure~3 assignments.
\end{enumerate}
\begin{verifybox}
T9: Hubbard model ($t=1$, $U=3$), exact diagonalisation in every $S_z$ sector. The ground-state $S$ equals
$\tfrac12|n_A-n_B|$ for all 7 graphs: CBD 0, TMM 1, TME 0, MX 1, TriMB $\tfrac32$, TetraMB 0, $p$-xylylene 0. \textbf{PASS}.
\end{verifybox}

\subsection{Disjoint diradicals: the singlet is selected by spin polarisation (Kollmar--Staemmler)}
For disjoint NBMOs (CBD, TME) the first-order gap $2K_{ab}$ vanishes (\cref{sec:huckel}). The degeneracy is
lifted only at second order, through the spin-polarised configurations above. In D$_{4h}$ CBD, ``the
bonding $\pi$-electrons can partially localize on the same pair of atoms as the same-spin unpaired
electrons; however, in the triplet state such advantageous spin polarization is not possible'' (review,
p.~443). The numerical signature is a gap $\propto U^2$, with the singlet lower.
\begin{verifybox}
T9b: gap $E_T-E_S>0$ (singlet lower), scaling as $U^{1.97}$ (CBD) and $U^{1.95}$ (TME) for $U\to0$;
cf.\ linear scaling with the triplet lower for the nondisjoint TMM and MX. \textbf{PASS}.
\end{verifybox}
\begin{sourcebox}
Kollmar \& Staemmler, \emph{Theor.\ Chim.\ Acta} \textbf{48}, 223 (1978), doi:10.1007/BF00549021 (not accessed).
Ovchinnikov, \emph{Theor.\ Chim.\ Acta} \textbf{47}, 297 (1978), doi:10.1007/BF00549259: $S=(N_A-N_B)/2$ as quoted in
arXiv:1906.08544. Lieb 1989: statement verified.
\end{sourcebox}
\begin{caveatbox}
Ovchinnikov's rule and Lieb's theorem apply to \emph{alternant} (bipartite) $\pi$ systems described by
Hubbard/Heisenberg-type models. The review's own counterexamples (pyridinium analogue V, Figure~6;
$\sigma$--$\pi$ diradicals such as the didehydrotoluenes) lie outside these assumptions. This is consistent
with its warning that such rules ``inevitably'' have ``limited scope''.
\end{caveatbox}

\section{Orbital energies: Koopmans, Janak, and the singlet--triplet gap as a HOMO difference}\label{sec:koopmans}

\begin{reviewclaim}
``An alternative theoretical interpretation of the Hund's multiplicity rule has been proposed on the
basis of Janak theorem (or extended Koopmans' theorem) in excited state density functional theory (DFT)
for multiplets: at the Hartree--Fock level, the singlet--triplet gap (STG) is equal to the difference
between the energies of the triplet's and open-shell singlet's HOMOs. This explanation echoes Pearson's
maximum hardness principle, but lacks dynamic electron correlation.'' (p.~442--443, ref.~20 Liu \&
Langenaeker 2003)
\end{reviewclaim}

\subsection{Koopmans' theorem}
For a single determinant of spin orbitals $\{\chi_i\}$ write $[ij|kl]$ for spin-orbital integrals. Then
\begin{equation}
  E_N=\sum_i\langle i|\hat h|i\rangle+\tfrac12\sum_{i,j}\bigl([ii|jj]-[ij|ji]\bigr).
\end{equation}
\begin{enumerate}
  \item Remove spin orbital $k$ \emph{without changing the others} (frozen orbitals). Every term containing
        $k$ disappears:
        $E_N-E_{N-1}^{(k)}=\langle k|\hat h|k\rangle+\tfrac12\sum_{j}\bigl([kk|jj]-[kj|jk]\bigr)+\tfrac12\sum_{i}\bigl([ii|kk]-[ik|ki]\bigr)$.
        The $i=j=k$ term, counted in both sums, is $[kk|kk]-[kk|kk]=0$.
  \item The two sums are equal by the symmetry of the integrals, so $E_N-E_{N-1}^{(k)}=\langle k|\hat h|k\rangle+\sum_j\bigl([kk|jj]-[kj|jk]\bigr)$.
  \item The right-hand side is the diagonal Fock element, i.e.\ the canonical orbital energy $\varepsilon_k$. Hence
        \begin{equation}\label{eq:koopmans}
          \mathrm{IP}_k\equiv E_{N-1}^{(k)}-E_N=-\varepsilon_k\qquad(\text{frozen orbitals}).
        \end{equation}
\end{enumerate}
Relaxation of the remaining orbitals lowers $E_{N-1}$. Correlation acts in the other direction, so
\cref{eq:koopmans} is an approximation with partially cancelling errors.

\subsection{Janak's theorem}
In Kohn--Sham DFT with fractional occupations $n_i$:
$E=\sum_in_i\langle\psi_i|-\tfrac12\nabla^2|\psi_i\rangle+\int v\rho+E_H[\rho]+E_{xc}[\rho]$, with
$\rho=\sum_in_i|\psi_i|^2$.
\begin{enumerate}
  \item Differentiate with respect to $n_i$. The explicit dependence (through the prefactor $n_i$ and through
        $\rho$) gives $\langle\psi_i|-\tfrac12\nabla^2|\psi_i\rangle+\int(v+v_H+v_{xc})|\psi_i|^2=\langle\psi_i|\hat h_{KS}|\psi_i\rangle=\varepsilon_i$.
  \item The implicit dependence through $\partial\psi_j/\partial n_i$ gives
        \[\sum_jn_j\bigl[\langle\partial_{n_i}\psi_j|\hat h_{KS}|\psi_j\rangle+\text{c.c.}\bigr]=\sum_jn_j\varepsilon_j\,\partial_{n_i}\langle\psi_j|\psi_j\rangle=0,\]
        because each $\psi_j$ is an eigenfunction of $\hat h_{KS}$ and stays normalised.
  \item Hence $\partial E/\partial n_i=\varepsilon_i$, for any $E_{xc}$ [verified: Janak, \emph{Phys.\ Rev.\ B} \textbf{18}, 7165 (1978), abstract].
        Integrating, $E(n_i=1)-E(n_i=0)=\int_0^1\varepsilon_i(n)\dd n\approx\varepsilon_i(\tfrac12)$ (Slater's transition state).
\end{enumerate}

\subsection{The singlet--triplet gap as a HOMO-energy difference (frozen-orbital derivation)}
Use the $2e/2o$ states of \cref{sec:energies} and define, for each state $X\in\{T,\text{OSS}\}$, the ``HOMO
energy'' as minus the energy needed to remove the electron from $\pb$, leaving $\pa$ frozen (the
Koopmans/$\Delta$SCF definition, \cref{eq:koopmans}). The ion has energy $E_{\text{ion}}=h_a$ in both cases.
\begin{enumerate}
  \item $\varepsilon_b^{T}=E_T-E_{\text{ion}}=h_b+J_{ab}-K_{ab}$.
  \item $\varepsilon_b^{\text{OSS}}=E_{\text{OSS}}-E_{\text{ion}}=h_b+J_{ab}+K_{ab}$.
  \item Subtract:
        \begin{equation}\label{eq:STeps}
          \boxed{E_{\text{OSS}}-E_T=\varepsilon_{\text{HOMO}}^{\text{OSS}}-\varepsilon_{\text{HOMO}}^{T}=2K_{ab}.}
        \end{equation}
\end{enumerate}
\begin{resultbox}
Within frozen orbitals the singlet--triplet gap \emph{is} the difference of the HOMO energies. The
triplet's HOMO lies lower by $2K_{ab}$: it is more strongly bound, i.e.\ ``harder''. That is the sense in
which the argument ``echoes'' the maximum-hardness principle. The sign convention is now fixed:
STG $\equiv E_S-E_T=\varepsilon^{S}_{\text{HOMO}}-\varepsilon^{T}_{\text{HOMO}}$.
Like Slater's argument, \cref{eq:STeps} has no dynamic correlation and no orbital relaxation, so it
inherits the limitations discussed in \cref{sec:hund}.
\end{resultbox}
\begin{verifybox}
T10 (\texttt{t10\_t12\_t13\_symbolic.py}): \cref{eq:STeps} symbolically. T2: the Koopmans identity
$\varepsilon_k=E_N-E_{N-1}^{(k)}$ checked numerically for random integrals in a 3-electron case. \textbf{PASS}.
\end{verifybox}
\begin{sourcebox}
Liu \& Langenaeker, \emph{Theor.\ Chem.\ Acc.}\ \textbf{110}, 338 (2003), doi:10.1007/s00214-003-0487-6: not accessed
(search summaries only state $\Delta E_{ST}=\Delta\varepsilon_{\text{HOMO}}$). The maximum-hardness connection is the
review's wording and was not verified against the original.
\end{sourcebox}

\section{Quasi-closed-shell (non-aufbau) radicals}\label{sec:qcs}

\begin{reviewclaim}
``In molecular systems with unpaired electrons, the violation of the aufbau principle in electron
configuration is reflected in the conversion of the SOMO and HOMO orbital energies, yielding a so-called
quasi-closed-shell state (Figure 10). Its main fingerprint is the electron attachment or removal pattern:
upon ionization an electron is removed from the higher energy doubly occupied MO, leaving another
unpaired electron behind, while upon addition new electron enters the former LUMO; essentially, SOMO does
not participate in these processes.'' (p.~448--449) ``\dots a rigorous support for the non-aufbau
occupation can be obtained a posteriori from the ionization potentials and electron affinities \dots as
well as from the determination of the spin state \dots of the corresponding open-shell
oxidation/reduction products.'' (p.~449)
\end{reviewclaim}

\subsection{Frozen-orbital model of the redox fingerprint}
Three relevant orbitals: SOMO $s$, the highest doubly occupied orbital $h$, and the LUMO $l$. Ground
configuration $(\text{core})\,s^1h^2$ (a doublet).
\begin{enumerate}
  \item UHF orbital energies, from \cref{eq:koopmans} applied to each spin orbital (Slater--Condon, \cref{sec:koopmans}):
        \begin{align}
          \varepsilon_h^\alpha&=h_h+J_{hh}+J_{hs}-K_{hs}, &\varepsilon_h^\beta&=h_h+J_{hh}+J_{hs},\label{eq:uhfh}\\
          \varepsilon_s^\alpha&=h_s+2J_{hs}-K_{hs}, &\varepsilon_s^\beta(\text{virtual})&=h_s+J_{ss}+2J_{hs}-K_{hs}.\label{eq:uhfs}
        \end{align}
        (Interactions with the core are absorbed into $h_h,h_s$.)
  \item \textbf{Spin polarisation of the donor HOMO (Scheme~3).} From \cref{eq:uhfh},
        $\varepsilon_h^\beta-\varepsilon_h^\alpha=K_{hs}\ge0$: the $\alpha$ electron of $h$ (parallel to the unpaired spin) is
        stabilised by exchange with the SOMO electron, while the $\beta$ electron is not. The donor's $\alpha$ and
        $\beta$ HOMOs split. That is the content of Scheme~3.
  \item \textbf{Condition for SOMO--HOMO conversion} (within the $\alpha$ manifold):
        $\varepsilon_s^\alpha<\varepsilon_h^\alpha\iff h_s+2J_{hs}-K_{hs}<h_h+J_{hh}+J_{hs}-K_{hs}\iff h_s+J_{hs}<h_h+J_{hh}$.
        In words: the SOMO's one-electron energy plus its repulsion with \emph{one} $h$ electron lies below
        $h$'s one-electron energy plus its \emph{self}-repulsion. A stable, low-lying radical orbital and a
        separate, electron-rich (high $h_h$, e.g.\ anionic) fragment favour this.
  \item \textbf{Oxidation.} Koopmans: $\mathrm{IP}(h)=-\varepsilon_h<\mathrm{IP}(s)=-\varepsilon_s$, so the electron leaves $h$.
        The product $(\text{core})\,s^1h^1$ is a diradical cation (or diradical, if the parent was an anion). Its
        triplet and open-shell singlet are split by $E_{\text{OSS}}-E_T=2K_{sh}$ (\cref{eq:EOSS} with $a,b\to s,h$).
        With $K_{sh}>0$ the triplet is lower: the observed ``triplet cation diradicals with nondegenerate SOMOs''.
  \item \textbf{Reduction.} The added electron can enter the $\beta$ SOMO, at cost $\varepsilon_s^\beta$ of \cref{eq:uhfs}, giving
        a closed-shell anion; or the LUMO, at $\varepsilon_l$, giving the diradical $s^1h^2l^1$. Koopmans picks the
        lower: the diradical forms iff $\varepsilon_l<\varepsilon_s^\beta=h_s+J_{ss}+2J_{hs}-K_{hs}$. The large on-site
        repulsion $J_{ss}$ makes this possible. Note that this is a \emph{separate} condition, not implied by
        step 3; the review's Figure~10 depicts the case where it holds, and Figure~17 gives examples.
  \item \textbf{Protonation switch.} Protonating the anionic fragment strongly lowers $h_h$ (Coulomb attraction of
        $\mathrm H^+$). Once $h_h+J_{hh}<h_s+J_{hs}$ the aufbau order is restored, and oxidation removes the SOMO
        electron, giving a closed-shell cation. This is the pH-switchable conversion of Figure~16.
\end{enumerate}

\begin{caveatbox}
Steps 1--5 use frozen orbitals and canonical orbital energies. By \cref{sec:rohf} orbital energies are
convention-dependent, and by \cref{sec:koopmans} relaxation shifts them. As the review says, the
\emph{observable} criterion is energetic: compare the two possible oxidised states directly,
$E(\text{diradical},T)$ vs.\ $E(\text{closed shell})$ ($\Delta$SCF).
\end{caveatbox}

% numbers from verification/logs/t11.log (T11b) and logs/t11c.log (T11c)
\subsection{Numerical test: a distonic radical anion and its conjugate acid}
Model: the aminoxyl radical $^\bullet$ON(Me)CH$_2$CH$_2$COO$^-$, with a carboxylate that is not conjugated to
the radical, and its protonated form $^\bullet$ON(Me)CH$_2$CH$_2$COOH. UB3LYP/6-31+G*; radical geometries
optimised (geomeTRIC). Vertical oxidation at the radical geometry to (i) the triplet (one electron removed
from below the SOMO, the diradical) and (ii) the closed-shell singlet (SOMO electron removed).
\begin{center}
\begin{tabular}{lcccc}
\toprule
radical & $\alpha$-SOMO position & IE$\to$triplet & IE$\to$closed shell & $\Delta$(T$-$CS)\\
\midrule
acid (COOH) & is the $\alpha$-HOMO & 9.489 eV & 8.179 eV & $+1.31$ eV\\
anion (COO$^-$) & 3 levels below the $\alpha$-HOMO & 4.008 eV & 3.946 eV & $+0.06$ eV\\
\bottomrule
\end{tabular}
\end{center}
\begin{enumerate}
  \item \emph{Orbital level: confirmed.} In the anion the $\alpha$-SOMO ($-2.40$~eV) lies below three doubly
        occupied levels (character not analysed) ($\alpha$-HOMO $-1.62$~eV). In the acid the SOMO is the HOMO. This is
        the pH-switchable SOMO--HOMO conversion of step 6.
  \item \emph{Energetic level: shifted, but no crossover.} Deprotonation moves the
        triplet--closed-shell preference by $1.25$~eV towards the diradical, but at this level the closed-shell
        product of the anion remains lower, by $0.06$~eV. For this short-linker model, the orbital inversion
        alone does not guarantee that oxidation yields the diradical. This is exactly why the review insists on
        the energetic, observable criterion (caveat above). T11c below shows that this verdict depends on the
        functional.
\end{enumerate}
% T11c results (functional and separation dependence)
\emph{(T11c pending.)}



\begin{verifybox}
T2: \cref{eq:uhfh,eq:uhfs}, $\varepsilon_h^\beta-\varepsilon_h^\alpha=K_{hs}$, and the cation $E_{\text{OSS}}-E_T=2K_{hs}$ verified
numerically against Slater--Condon determinant energies for random integrals. \textbf{PASS}.
\end{verifybox}
\begin{sourcebox}
Gryn'ova, Marshall, Blanksby \& Coote, \emph{Nat.\ Chem.}\ \textbf{5}, 474 (2013); Gryn'ova \& Coote,
\emph{J.\ Am.\ Chem.\ Soc.}\ \textbf{135}, 15392 (2013): summaries only (SOMO--HOMO conversion when a stabilised
radical is not $\pi$-conjugated with an anionic group; protonation restores normal ordering). Westcott et al.\
(2000): VO(oep), where the unpaired metal d electron is bound more strongly than several ligand pairs
(author summary).
\end{sourcebox}

\section{ROHF orbital energies are not unique (canonicalisation)}\label{sec:rohf}

\begin{reviewclaim}
``\dots the orbitals themselves are neither observable, nor unique (in fact, erroneous non-aufbau
occupations can be observed with restricted open-shell methods depending on the canonicalization
parameters of the $\alpha$ and $\beta$ electrons in Fock matrix, see Ref 82) \dots'' (p.~449; ref.~82
Glaesemann \& Schmidt 2010)
\end{reviewclaim}

\subsection{High-spin ROHF and its stationarity conditions}
Orbital spaces: doubly occupied (core) $c$, singly occupied ($\alpha$ only) $o$, virtual $v$. Densities:
$\mathbf D^\alpha=\mathbf D_c+\mathbf D_o$ and $\mathbf D^\beta=\mathbf D_c$. The energy has the UHF form
$E[\mathbf D^\alpha,\mathbf D^\beta]$, so its first variation is
\begin{equation}\label{eq:dE}
  \delta E=\operatorname{tr}(\mathbf F^\alpha\,\delta\mathbf D^\alpha)+\operatorname{tr}(\mathbf F^\beta\,\delta\mathbf D^\beta),
\end{equation}
with $\mathbf F^\sigma=\partial E/\partial\mathbf D^\sigma$ the UHF Fock matrices built from the ROHF densities.
\begin{enumerate}
  \item An infinitesimal real rotation mixing orbital $q$ into orbital $p$ by angle $\kappa$ changes the density of a
        spin $\sigma$ only if exactly one of $p,q$ is occupied in that spin. Then
        $\delta\mathbf D^\sigma=\kappa(\ket p\bra q+\ket q\bra p)$, and $\delta E$ picks up $2\kappa F^\sigma_{pq}$.
  \item \textbf{$c$--$o$ rotation:} for $\alpha$ both orbitals are occupied (no change); for $\beta$, $c$ is occupied
        and $o$ empty. Stationarity requires $F^\beta_{co}=0$.
  \item \textbf{$o$--$v$ rotation:} for $\alpha$, $o$ is occupied and $v$ empty; for $\beta$ both are empty.
        Stationarity requires $F^\alpha_{ov}=0$.
  \item \textbf{$c$--$v$ rotation:} both spins change. Stationarity requires $F^\alpha_{cv}+F^\beta_{cv}=0$.
  \item \textbf{Rotations within $c$, within $o$, or within $v$} leave $\mathbf D^\alpha$ and $\mathbf D^\beta$ unchanged, hence
        $\delta E=0$ \emph{identically}. Nothing in the variational principle fixes them.
\end{enumerate}

\subsection{Effective Fock operator and the arbitrariness of orbital energies}
To get orbitals from a single eigenvalue problem, one builds an effective Fock matrix $\mathbf R$ whose
off-diagonal blocks enforce steps 2--4:
\begin{equation}
  \mathbf R=\begin{pmatrix}
  \mathbf R_{cc} & \mathbf F^\beta_{co} & \tfrac12(\mathbf F^\alpha+\mathbf F^\beta)_{cv}\\
  \cdot & \mathbf R_{oo} & \mathbf F^\alpha_{ov}\\
  \cdot & \cdot & \mathbf R_{vv}\end{pmatrix},\qquad
  \mathbf R_{xx}=A_{xx}\mathbf F^\alpha_{xx}+B_{xx}\mathbf F^\beta_{xx}\ \ (\text{any choice}).
\end{equation}
At convergence the off-diagonal blocks vanish (steps 2--4), and the orbitals diagonalise each diagonal
block separately. By step 5 the total energy, the densities and every observable derived from them are
independent of $A_{xx},B_{xx}$. The \emph{orbital energies}, the eigenvalues of $\mathbf R_{xx}$, are not. The
relative position of the SOMO (eigenvalue of $\mathbf R_{oo}$) and the HOMO of the closed shell (top
eigenvalue of $\mathbf R_{cc}$) is therefore a matter of convention.

\begin{verifybox}
T11a (\texttt{t11\_rohf\_and\_distonic.py}), ROHF/6-31G*:
\begin{itemize}
  \item Stationarity: $|F^\beta_{co}|,|F^\alpha_{ov}|,|(F^\alpha+F^\beta)_{cv}|<10^{-7}$ for H$_2$NO, NO$_2$, CH$_3$ and allyl.
        The blocks that need not vanish do not: $|F^\alpha_{co}|=0.11$, $|F^\beta_{ov}|=0.16$ (H$_2$NO). This confirms
        steps 2--4. PySCF's \texttt{get\_roothaan\_fock} uses exactly these off-diagonal blocks.
  \item Canonicalisation, H$_2$NO: $\varepsilon_{\text{SOMO}}=-0.119$ ($\mathbf R_{oo}=\tfrac12(\mathbf F^\alpha+\mathbf F^\beta)$),
        $-0.440$ ($\mathbf F^\alpha$), $+0.203\,E_h$ ($\mathbf F^\beta$), while the closed-shell HOMO is $-0.467$. Same energy
        $E=-130.36970936\,E_h$ in every case.
  \item With one open orbital, $\varepsilon_{\text{SOMO}}(A)=AF^\alpha_{oo}+(1-A)F^\beta_{oo}$ is linear in $A$. The SOMO drops
        below the HOMO for $A>A^*$, with $A^*\approx1.04$ for H$_2$NO and NO$_2$. The flip was verified explicitly.
\end{itemize}
\textbf{PASS}.
\end{verifybox}
\begin{caveatbox}
For these ordinary radicals the ordering flips only for weights slightly outside the ``convex'' range
$0\le A\le1$. So the review's warning is mathematically exact (the ordering is not invariant), but whether
a \emph{standard} canonicalisation produces a spurious non-aufbau ordering depends on the molecule and
the scheme. This is why the review, correctly, asks for evidence from observables: ionisation
energies, redox potentials and the spin state of the products (\cref{sec:qcs}).
\end{caveatbox}
\begin{sourcebox}
Glaesemann \& Schmidt, \emph{J.\ Phys.\ Chem.\ A} \textbf{114}, 8772 (2010), doi:10.1021/jp101758y: not accessed.
The specific canonicalisation parameter sets they compare could not be checked. The derivation above
is general.
\end{sourcebox}

\section{The 3d/4s problem}\label{sec:3d4s}

\begin{reviewclaim}
``In neutral atoms, the energy of a 4s orbital is slightly above that of a 3d orbital (when both orbitals
are populated), therefore, 4s electrons are first ionized. However, 4s orbital is also populated first,
because it offers less repulsion with a 3d electron than that between two 3d electrons, once the orbital
energies are re-evaluated in the new configurations.'' (p.~448, refs.~79--81)
\end{reviewclaim}

\subsection{Average-of-configuration energies}
Model a configuration $d^{n_d}s^{n_s}$ by its average energy. Each electron has one-electron energy $h_d$ or $h_s$,
and every pair of electrons repels with an \emph{average} repulsion that depends only on the shells:
$\bar J_{dd}$ (two d electrons), $\bar J_{ds}$, and $J_{ss}$. (For example $\bar J_{dd}=F^0(dd)-\tfrac{2}{63}[F^2(dd)+F^4(dd)]$
in Slater--Condon parameters. Only the shell structure matters here.)
\begin{enumerate}
  \item Count the pairs: $\binom{n_d}2$ d--d pairs, $n_dn_s$ d--s pairs, $\binom{n_s}2$ s--s pairs. Hence
        \begin{equation}
          E(d^{n_d}s^{n_s})=n_dh_d+n_sh_s+\tbinom{n_d}{2}\bar J_{dd}+n_dn_s\bar J_{ds}+\tbinom{n_s}{2}J_{ss}.
        \end{equation}
  \item The two competing configurations with $n+2$ valence electrons:
        \begin{align*}
          E(d^ns^2)&=nh_d+2h_s+\tbinom n2\bar J_{dd}+2n\bar J_{ds}+J_{ss},\\
          E(d^{n+1}s^1)&=(n+1)h_d+h_s+\tbinom{n+1}2\bar J_{dd}+(n+1)\bar J_{ds}.
        \end{align*}
  \item Subtract, using $\binom{n+1}2-\binom n2=n$:
        \begin{equation}\label{eq:Delta3d4s}
          \Delta\equiv E(d^{n+1}s)-E(d^ns^2)=h_d-h_s+n\bar J_{dd}+(1-n)\bar J_{ds}-J_{ss}.
        \end{equation}
  \item Orbital energies \emph{in the configuration $d^ns^2$}, Koopmans style ($\varepsilon=E(N)-E(N-1)$, frozen):
        \begin{align*}
          \varepsilon_d&=E(d^ns^2)-E(d^{n-1}s^2)=h_d+(n-1)\bar J_{dd}+2\bar J_{ds},\\
          \varepsilon_s&=E(d^ns^2)-E(d^ns^1)=h_s+n\bar J_{ds}+J_{ss}.
        \end{align*}
  \item Therefore $\varepsilon_d-\varepsilon_s=h_d-h_s+(n-1)\bar J_{dd}+(2-n)\bar J_{ds}-J_{ss}$, and comparing with \cref{eq:Delta3d4s}:
        \begin{equation}\label{eq:3d4s}
          \boxed{\Delta=(\varepsilon_d-\varepsilon_s)+\bigl(\bar J_{dd}-\bar J_{ds}\bigr).}
        \end{equation}
\end{enumerate}
\begin{resultbox}
Even when $\varepsilon_{3d}<\varepsilon_{4s}$ (3d ``below'' 4s), the $d^ns^2$ configuration is lower in energy
($\Delta>0$) whenever $\bar J_{dd}-\bar J_{ds}>\varepsilon_s-\varepsilon_d$. Because 3d is compact and 4s diffuse,
$\bar J_{dd}\gg\bar J_{ds}$: moving an electron from 4s into 3d costs the large d--d repulsion. Example
(illustrative parameters, $n=2$): $\varepsilon_d=-1.50$, $\varepsilon_s=-0.90$, $\Delta=+0.10$.
\end{resultbox}
\paragraph{Ionisation.} From $d^ns^2$, the Koopmans ionisation energies are $-\varepsilon_s$ and $-\varepsilon_d$.
With $\varepsilon_s>\varepsilon_d$, the 4s electron leaves first, as observed for transition-metal cations.
Both statements in the review follow from the same two numbers, $\varepsilon$ and $\bar J$, evaluated
\emph{in the relevant configuration}. The ``building-up'' order and the ``removal'' order need not agree.
\begin{verifybox}
T12 (\texttt{t10\_t12\_t13\_symbolic.py}): \cref{eq:Delta3d4s,eq:3d4s} and both orbital-energy formulas
verified symbolically for general $n$; the cation analogue is $\Delta(n-1)$; the numerical example satisfies
$\varepsilon_d<\varepsilon_s$ and $\Delta>0$. \textbf{PASS}.
\end{verifybox}
\begin{sourcebox}
Vanquickenborne, Pierloot \& Devoghel, \emph{Inorg.\ Chem.}\ \textbf{28}, 1805 (1989); Melrose \& Scerri,
\emph{J.\ Chem.\ Educ.}\ \textbf{73}, 498 (1996): not accessed. The derivation reconstructs the
average-of-configuration argument independently.
\end{sourcebox}

\section{Why individual electrons have no good quantum numbers}\label{sec:qn}

\begin{reviewclaim}
``\dots according to modern quantum mechanics, in the stationary states the quantum numbers of individual
electrons do not commute with the Hamiltonian, only the total angular momentum does. For this reason
\dots\ the Pauli exclusion principle has been reformulated to reflect the antisymmetry of the
wavefunction \dots'' (p.~448, ref.~77 Scerri 1991). ``\dots assigning individual electrons to individual
orbitals and individual sets of quantum numbers is erroneous'' (p.~441).
\end{reviewclaim}

\subsection{The commutators}
Take an atom (nucleus at the origin), so $\Hop=\sum_i\hat h(i)+\sum_{i<j}r_{ij}^{-1}$ with
$\hat h(i)=-\tfrac12\nabla_i^2-Z/r_i$. Orbital angular momentum of electron 1:
$\hat{\mathbf l}_1=-i\,\mathbf r_1\times\nabla_1$.
\begin{enumerate}
  \item $\hat{\mathbf l}_1$ commutes with $\hat h(1)$ (spherical symmetry) and with $\hat h(j)$, $j\ne1$ (different variables).
        The only possible obstruction is $1/r_{12}$.
  \item Apply $\hat l_{1z}=-i(x_1\partial_{y_1}-y_1\partial_{x_1})$ to $V=1/r_{12}$, with
        $\partial_{y_1}r_{12}^{-1}=-(y_1-y_2)/r_{12}^3$ and $\partial_{x_1}r_{12}^{-1}=-(x_1-x_2)/r_{12}^3$:
        \[
          \hat l_{1z}V=-i\,\frac{-x_1(y_1-y_2)+y_1(x_1-x_2)}{r_{12}^3}=-i\,\frac{x_1y_2-x_2y_1}{r_{12}^3}.
        \]
        Since $\hat l_{1z}$ is a first-order differential operator, $[\hat l_{1z},V]f=(\hat l_{1z}V)f$ for any $f$, so
        \begin{equation}
          [\hat l_{1z},\,r_{12}^{-1}]=-i\,\frac{x_1y_2-x_2y_1}{r_{12}^3}\neq0 .
        \end{equation}
  \item For electron 2, $\partial_{y_2}r_{12}^{-1}=+(y_1-y_2)/r_{12}^3$ etc., which gives
        $\hat l_{2z}V=-i\,(x_2y_1-x_1y_2)/r_{12}^3=-\hat l_{1z}V$.
  \item Hence $[\hat L_z,r_{12}^{-1}]=[\hat l_{1z}+\hat l_{2z},r_{12}^{-1}]=0$. The same holds for the $x$ and $y$
        components. Physically, $\hat{\mathbf L}$ generates a simultaneous rotation of both electrons, which
        leaves $r_{12}$ unchanged; $\hat{\mathbf l}_1$ rotates electron 1 alone, which changes $r_{12}$.
  \item Consequently $[\Hop,\hat{\mathbf L}]=0$ and $[\Hop,\hat L^2]=0$, but $[\Hop,\hat{\mathbf l}_1]\neq0$ and $[\Hop,\hat l_1^2]\neq0$.
        $\hat{\mathbf S}$ and $\hat S^2$ commute with $\Hop$ trivially, since $\Hop$ is spin-free.
\end{enumerate}
\begin{resultbox}
Exact atomic eigenstates can be labelled by $L,M_L,S,M_S$ (and parity), not by one-electron $(n,l,m_l,m_s)$.
Those labels belong to the independent-particle model, in which $1/r_{ij}$ is replaced by a spherically
averaged mean field. Pauli's 1925 one-electron formulation is therefore a statement about that model. The
exact, model-independent statement is the antisymmetry \cref{eq:antisym}.
\end{resultbox}
\begin{verifybox}
T13 (\texttt{t10\_t12\_t13\_symbolic.py}): $[\hat l_{1z},r_{12}^{-1}]f$ evaluated symbolically on an arbitrary
function $f$ reproduces the expression above. $[\hat L_k,r_{12}^{-1}]f=0$ for $k=x,y,z$, and $[\hat l_1^2,r_{12}^{-1}]\neq0$. \textbf{PASS}.
\end{verifybox}
\begin{sourcebox}
Scerri, \emph{Br.\ J.\ Philos.\ Sci.}\ \textbf{42}, 309 (1991); Pauli, \emph{Z.\ Phys.}\ \textbf{31}, 765 (1925): not
accessed. The commutator algebra is standard and was verified here.
\end{sourcebox}

\section{The adiabatic principle and the aufbau argument}\label{sec:adiabatic}

\begin{reviewclaim}
Bohr's building-up principle assumed ``that electrons exist in sharp stationary states, and that these
states are preserved upon addition of an electron and a proton''. These assumptions ``originate from the
Ehrenfest's adiabatic principle: `if a system be affected in a reversible adiabatic way, the allowed
motions are transformed into allowed motions.' However, this principle does not apply to aperiodic
systems, i.e., all systems other than a hydrogen atom.'' ``\dots the aufbau principle, by assuming the
existence of unchanging stationary individual electronic states, violates the Pauli exclusion principle.''
(p.~448)
\end{reviewclaim}

\subsection{The quantum adiabatic theorem (what it does guarantee)}
Let $\Hop(t)$ change slowly, with instantaneous eigenstates $\Hop(t)\psi_n(t)=E_n(t)\psi_n(t)$, non-degenerate.
\begin{enumerate}
  \item Expand $\Psi(t)=\sum_nc_n(t)e^{i\theta_n(t)}\psi_n(t)$ with $\theta_n=-\int_0^tE_n\dd t'$.
  \item Insert into $i\partial_t\Psi=\Hop\Psi$. The $E_n$ terms cancel against $\dot\theta_n$, leaving
        $\sum_n\bigl(\dot c_n\psi_n+c_n\dot\psi_n\bigr)e^{i\theta_n}=0$.
  \item Project on $\psi_m$: $\dot c_m=-\sum_nc_n\braket{\psi_m}{\dot\psi_n}e^{i(\theta_n-\theta_m)}$.
  \item Differentiate $\Hop\psi_n=E_n\psi_n$ in time and project on $\psi_m$ ($m\ne n$):
        $\braket{\psi_m}{\dot\psi_n}=\dfrac{\mel{\psi_m}{\dot{\Hop}}{\psi_n}}{E_n-E_m}$.
  \item If $|\mel{\psi_m}{\dot\Hop}{\psi_n}|\ll(E_n-E_m)^2$, the rapidly oscillating phase makes the coupling average
        out, and $|c_m|$ stays at its initial value: a system starting in $\psi_n$ ends in the
        \emph{continuously connected} eigenstate $\psi_n(t)$.
\end{enumerate}

\subsection{What it does not guarantee: the aufbau assumptions}
\begin{enumerate}
  \item The theorem is about a \emph{fixed} Hilbert space and a Hamiltonian with a continuously varying
        parameter. Going from element $Z$ to $Z+1$ adds a particle. The $N$- and $(N+1)$-electron Hilbert
        spaces are different, so there is no continuous path of Hamiltonians on which the theorem could act.
  \item Even if one varies $Z$ continuously at fixed $N$, the theorem preserves the label of an
        \emph{$N$-electron} eigenstate, the conserved quantities of \cref{sec:qn} ($L,S,\dots$). It says nothing about
        one-electron states. These are not eigenstates of $\Hop$, because $[\Hop,\hat{\mathbf l}_i]\ne0$.
  \item Within the orbital model, the orbitals of $d^ns^2$ and $d^{n+1}s$ differ (\cref{sec:3d4s}). Adding or removing
        an electron changes all orbitals (\cref{sec:koopmans}, orbital relaxation beyond Koopmans).
\end{enumerate}
\begin{resultbox}
The adiabatic theorem cannot justify the idea that previously ``built-up'' one-electron states persist,
unchanged, when electrons and protons are added. That is the mathematical content behind the review's
historical criticism. The claim that Ehrenfest's principle ``does not apply to aperiodic systems'' refers
to the old (Bohr--Sommerfeld) quantum theory, where adiabatic invariants $\oint p\dd q$ are defined only
for (multiply) periodic motion. It is historical and could not be checked against the original here.
\end{resultbox}
\begin{sourcebox}
Ehrenfest, \emph{Philos.\ Mag.}\ \textbf{33}, 500 (1917); Scerri (1991): not accessed (quotation from the review).
Born \& Fock (1928) for the quantum adiabatic theorem; the derivation above is standard.
\end{sourcebox}

\section{Broken-symmetry DFT/UHF and spin projection}\label{sec:bs}

\begin{reviewclaim}
``In terms of computations, relatively `cheap' methods such as unrestricted or broken symmetry formalisms
suffer from a number of serious drawbacks \dots'' (p.~443)
\end{reviewclaim}

\subsection{Why a single unrestricted determinant cannot be an open-shell singlet}
Take the $2e/2o$ broken-symmetry (BS) determinant $\ket{a'\bar b'}$ with $a'$ localised mostly on centre $l$ and
$b'$ mostly on $r$. In the orthonormal basis $\pa,\pb$ (overlap $\braket{a'}{b'}=s$) it is a mixture of the
$M_S=0$ singlet and triplet (\cref{sec:twoe}, \cref{eq:T0,eq:OSS}). Its spin expectation value is
\begin{equation}
  \ev{\hat S^2}_{\text{BS}}=S_z(S_z+1)+N_\beta-\sum_{i\in\alpha,\,j\in\beta}|\braket{i}{j}|^2=0+1-s^2,
\end{equation}
which lies between $0$ (pure singlet, $s=1$) and $1$ (equal singlet--triplet mixture, $s=0$).

\subsection{Yamaguchi's approximate projection}
\begin{enumerate}
  \item Model: $\hat H=-2J\,\hat{\mathbf S}_1\cdot\hat{\mathbf S}_2$. With $\hat{\mathbf S}_1\cdot\hat{\mathbf S}_2=\tfrac12[S(S+1)-\tfrac32]$:
        $E_T=-\tfrac J2$, $E_S=\tfrac{3J}2$, so $E_S-E_T=2J$.
  \item Write the BS state as $\Psi_{\text{BS}}=c_S\Psi_S+c_T\Psi_T$ ($M_S=0$), so $\ev{\hat S^2}_{\text{BS}}=2|c_T|^2$ and
        $E_{\text{BS}}=|c_S|^2E_S+|c_T|^2E_T$.
  \item Then $E_{\text{BS}}-E_T=|c_S|^2(E_S-E_T)=\bigl(1-\tfrac12\ev{\hat S^2}_{\text{BS}}\bigr)2J=\bigl(2-\ev{\hat S^2}_{\text{BS}}\bigr)J$.
  \item With $E_{\text{HS}}=E_T$ (single determinant, exact in this model) and $\ev{\hat S^2}_{\text{HS}}=2$:
        \begin{equation}\label{eq:yamaguchi}
          \boxed{J=\frac{E_{\text{BS}}-E_{\text{HS}}}{\ev{\hat S^2}_{\text{HS}}-\ev{\hat S^2}_{\text{BS}}}},\qquad E_S-E_T=2J .
        \end{equation}
\end{enumerate}
\begin{verifybox}
T14 (\texttt{t06\_t09\_models.py}): BS-UHF solved self-consistently for the Hubbard dimer.
Yamaguchi's $J$ compared with the exact $J=(E_S-E_T)/2=\tfrac14(U-\sqrt{U^2+16t^2})$:
relative error $-7.7\%$ ($U=3t$), $-0.9\%$ ($6t$), $-0.1\%$ ($12t$), $<0.05\%$ ($24t$). The projection works when the
two centres are weakly coupled (strongly correlated, $\ev{\hat S^2}_{\text{BS}}\to1$) and degrades as the
system becomes more covalent. This is one of the ``serious drawbacks'' of the cheap approaches. \textbf{PASS}.
\end{verifybox}
\begin{sourcebox}
Formula as given by secondary sources (IntechOpen chapter, ORCA manual) and attributed to Yamaguchi and
co-workers (1986/1988). Original papers not accessed. The derivation above is self-contained. Sign
conventions differ between sources; the one used here is $\hat H=-2J\hat{\mathbf S}_1\cdot\hat{\mathbf S}_2$, with $J<0$
antiferromagnetic.
\end{sourcebox}

\section{Summary: verification record, corrections and sources}\label{sec:summary}

\subsection{Verification record}
All scripts are in \texttt{verification/} and run with
\texttt{make verify} / \texttt{uv run --project . python -I <script>} (uv-locked: PySCF 2.14, SymPy 1.14, NumPy 2.5). Logs are in
\texttt{plans/tests.md}.
\begin{center}
\small
\begin{tabular}{llll}
\toprule
test & what is checked & section & result\\
\midrule
T1 & spin eigenfunctions; Fig.~1 combinations; sign erratum & \ref{sec:anti}, \ref{sec:twoe} & PASS\\
T2 & $2e/2o$ singlet matrix and $E_T$ vs.\ full CI; Koopmans relations & \ref{sec:energies}, \ref{sec:koopmans}, \ref{sec:qcs} & PASS\\
T3 & $K_{pq}\ge0$, $J_{pq}\ge K_{pq}$ (water MOs) & \ref{sec:holes} & PASS\\
T4 & He $2\,^3S/2\,^1S$: virial, less screening & \ref{sec:hund} & PASS\\
T5 & H$_2$ minimal basis: RHF/CI/triplet/OSS limits & \ref{sec:static} & PASS\\
T6 & two-orbital gap model & \ref{sec:static} & PASS\\
T7 & extended Hubbard dimer, kinetic exchange & \ref{sec:kinex} & PASS\\
T8 & Hückel NBMO rules, disjointness & \ref{sec:huckel} & PASS\\
T9 & Hubbard ED: Ovchinnikov/Lieb spins; spin alternation & \ref{sec:spinpol} & PASS\\
T9b & gap $\propto U$ (nondisjoint) vs.\ $\propto U^2$ (disjoint) & \ref{sec:huckel}, \ref{sec:spinpol} & PASS\\
T10 & S--T gap $=$ HOMO-energy difference & \ref{sec:koopmans} & PASS\\
T11a & ROHF stationarity, canonicalisation dependence & \ref{sec:rohf} & PASS\\
T11b & distonic anion (B3LYP): orbital inversion; no $\Delta$SCF crossover & \ref{sec:qcs} & FAIL (diagnosed by T11c)\\
T11c & CAM-B3LYP: diradical lower by 0.36~eV; acid closed shell; pH switch & \ref{sec:qcs} & PASS (partial run)\\
T12 & 3d/4s algebra & \ref{sec:3d4s} & PASS\\
T13 & $[\hat l_1,r_{12}^{-1}]\ne0$, $[\hat L,r_{12}^{-1}]=0$ & \ref{sec:qn} & PASS\\
T14 & Yamaguchi projection on BS-UHF dimer & \ref{sec:bs} & PASS\\
T15 & same-spin vs.\ opposite-spin correlation & \ref{sec:hund} & PASS (claim qualified)\\
T17 & Lean~4/Mathlib proofs of the algebraic identities & several & PASS (17 thms)\\
\bottomrule
\end{tabular}
\end{center}

\paragraph{Lean formalisation (T17).} \texttt{verification/lean\_proofs/Ch5103/Identities.lean} contains
machine-checked proofs of: the factorisations \cref{eq:sum,eq:diff} and a counterexample showing the
review's ``$+$'' combination is not the singlet; $E_S-E_T=2K$ as a HOMO-energy difference; the scaling
derivative $\dd E/\dd\eta|_1=2T+V$ and its virial consequences; the less-screening inequalities
\cref{eq:lessscreen}; the incompatibility of frozen orbitals with the virial theorem; the H$_2$ dissociation
limits; the gap-model and extended-dimer secular roots; the Hückel pairing theorem; the 3d/4s identity
\cref{eq:3d4s}; and Yamaguchi's formula \cref{eq:yamaguchi}.

\subsection{Corrections and qualifications to the review}
\begin{enumerate}
  \item \textbf{Sign error} (p.~441): the open-shell singlet is $\tfrac12(D_{\text{iv}}-D_{\text{iii}})$; the printed
        $\tfrac12(D_{\text{iii}}+D_{\text{iv}})$ is the $M_S=0$ triplet (\cref{sec:twoe}).
  \item \textbf{Citation year}: Katriel \& Pauncz, \emph{Adv.\ Quantum Chem.}\ \textbf{10}, 143--185 is 1977, not 1975.
  \item \textbf{Same-spin correlation} ``an order of magnitude'' smaller: true for He ($\sim$1/40), but in molecular
        valence shells at MP2 the per-pair ratio is $\sim$0.3--0.4 (\cref{sec:hund}).
  \item \textbf{Janak/Koopmans statement}: the sign convention is unstated in the review; with frozen orbitals
        $E_S-E_T=\varepsilon^S_{\text{HOMO}}-\varepsilon^T_{\text{HOMO}}=2K$ (\cref{sec:koopmans}).
  \item \textbf{Longuet-Higgins count} (used implicitly for NBMO counting): $N-2T$ is a lower bound
        (cyclobutadiene exceeds it).
  \item \textbf{ROHF canonicalisation}: the ordering is convention-dependent, as stated. For common radicals
        a flip needs weights slightly outside $[0,1]$, so observables must decide (\cref{sec:rohf}).
\end{enumerate}

\subsection{Sources and access status}
\begin{center}
\footnotesize
\begin{tabular}{p{0.52\textwidth}p{0.40\textwidth}}
\toprule
reference & status in this work\\
\midrule
Hongo et al., J.\ Chem.\ Phys.\ 121, 7144 (2004), doi:10.1063/1.1795151 & full text (arXiv cond-mat/0408147); Table I used\\
Boyd, Nature 310, 480 (1984), doi:10.1038/310480a0 & abstract\\
Katriel \& Pauncz, Adv.\ Quantum Chem.\ 10, 143 (1977), doi:10.1016/S0065-3276(08)60580-8 & bibliographic only\\
Krylov, Chem.\ Phys.\ Lett.\ 338, 375 (2001) & introduction (preview)\\
Balkov\'a \& Bartlett, J.\ Chem.\ Phys.\ 101, 8972 (1994), doi:10.1063/1.468057 & abstract\\
Lieb, Phys.\ Rev.\ Lett.\ 62, 1201 (1989) & theorem statement\\
Ovchinnikov, Theor.\ Chim.\ Acta 47, 297 (1978), doi:10.1007/BF00549259 & as quoted in arXiv:1906.08544\\
Janak, Phys.\ Rev.\ B 18, 7165 (1978) & abstract\\
Borden \& Davidson, J.\ Am.\ Chem.\ Soc.\ 99, 4587 (1977), doi:10.1021/ja00456a010 & not accessed\\
Kollmar \& Staemmler, Theor.\ Chim.\ Acta 48, 223 (1978), doi:10.1007/BF00549021 & not accessed\\
Liu \& Langenaeker, Theor.\ Chem.\ Acc.\ 110, 338 (2003), doi:10.1007/s00214-003-0487-6 & not accessed\\
Yamanaka et al., Int.\ J.\ Quantum Chem.\ 105, 687 (2005) & not accessed\\
Glaesemann \& Schmidt, J.\ Phys.\ Chem.\ A 114, 8772 (2010), doi:10.1021/jp101758y & not accessed\\
Vanquickenborne et al., Inorg.\ Chem.\ 28, 1805 (1989); Melrose \& Scerri (1996) & not accessed\\
Gryn'ova et al., Nat.\ Chem.\ 5, 474 (2013); Gryn'ova \& Coote, J.\ Am.\ Chem.\ Soc.\ 135, 15392 (2013) & abstracts/summaries\\
Westcott et al., J.\ Am.\ Chem.\ Soc.\ 122, 8083 (2000), doi:10.1021/ja994018p & bibliographic + author summary\\
Scerri (1991); Ehrenfest (1917); Pauli (1925); Slater (1929) & not accessed\\
\bottomrule
\end{tabular}
\end{center}
Where an original was not accessed, nothing in this document is attributed to it beyond what the review
itself states. The mathematics was derived independently and verified as listed above.


\end{document}
